% concatenated from main.tex, introduction.tex, random_maps.tex, convex_billiards.tex, random_additive_Birkhoff.tex, standard.tex
\documentclass[a4paper,12pt,reqno]{amsart}

% ---------- Language, encoding, fonts ----------
\usepackage[english]{babel}
\usepackage[T1]{fontenc}
\usepackage[utf8]{inputenc}

% Font: TeX Gyre Termes (text) + Termes Math (math)

\usepackage{tgpagella}
\usepackage{mathpazo}

%\usepackage{tgtermes}
%\usepackage{newtxmath}

% Micro-typography
\usepackage{microtype}

% ---------- Math ----------
\usepackage{mathtools}        % extends amsmath
\usepackage{amssymb}
\usepackage{amsthm}
\usepackage[mathscr]{eucal}
\usepackage{bbm}
% \usepackage{amscd}          % if you need commutative diagrams

% ---------- Graphics ----------
\usepackage{graphicx}
% \usepackage{epstopdf}       % only if you still use .eps files
\usepackage{subcaption}       % better than subfig

% ---------- Spacing ----------
\usepackage{setspace}
% \doublespacing              % enable if needed

% ---------- Colors & changes ----------
\usepackage[dvipsnames]{xcolor}
\usepackage[draft]{changes}   % use [final] for clean version

% ---------- Hyperlinks ----------
\usepackage[
  colorlinks = true,
  linkcolor  = blue,
  citecolor  = blue,
  urlcolor   = blue
]{hyperref}

%\usepackage{stmaryrd}
%\usepackage{standalone}
%\usepackage{refcheck}

%\usepackage[pagebackref]{hyperref}
%\renewcommand*{\backref}[1]{Cited at page [#1]}

%%%%%%%%%%%%%%%%%%%%%%%%%%%%%%%%%%%%%%%%%%%%%%%%%%%%%
%\addtolength{\evensidemargin}{-15mm}
%\addtolength{\oddsidemargin}{-15mm}
%\addtolength{\textwidth}{30mm}
%\addtolength{\textheight}{20mm}
%\addtolength{\topmargin}{-10mm}
%%%%%%%%%%%%%%%%%%%%%%%%%%%%%%%%%%%%%%%%%%%%%%%%%%%%
\renewcommand{\theequation}{\thesection.\arabic{equation}}
\numberwithin{equation}{section}
%%%%%%%%%%%%%%%%%%%%%%%%%%%%%%%%%%%%%%%%%%%%%%%%%%%%%
%\newtheorem{thm}{Theorem}[section]
%\newtheorem{lem}[thm]{Lemma}%[section]
%\newtheorem{prop}[thm]{Proposition}%[section]
%\newtheorem{cor}[thm]{Corollary}%[section]
%%\newtheorem{Jac}{Jacobi}

%%%%%%%%%%%%%%%%%%%%%%%%%%%%%%%%%%%%%%%%%%%%%%%%%%%%%
\newtheorem{theorem}{Theorem}[section]
\newtheorem{proposition}[theorem]{Proposition}
\newtheorem{conjecture}[theorem]{Conjecture}
\newtheorem{corollary}[theorem]{Corollary}
\newtheorem{lemma}[theorem]{Lemma}
\newtheorem{estimate}[theorem]{Estimate}
%\newtheorem{assumption}[theorem]{Assumption}

\newtheorem{assumption}{Assumption}
\renewcommand{\theassumption}{\Alph{assumption}}

{\theoremstyle{definition}%{plain}
{%\theorembodyfont{\normalfont\rmfamily}
%\newthbeorem{definition}[theorem]{Definition}
\newtheorem{remark}[theorem]{Remark}
\newtheorem{example}[theorem]{Example}
\newtheorem{exercise}[section]{Exercise}
\newtheorem{defn}[theorem]{Definition}
}}

%%%%%%%%%%%%%%%%%%%%%%%%%%%%%%%%%%%%%%%%%%%%%%%%%%%%


\newcommand{\cal}{\mathcal}

\newcommand{\A}{{\cal A}}
\newcommand{\BB}{{\cal B}}
\newcommand{\CC}{{\cal C}}
\newcommand{\CCC}{{\cal C}}
\newcommand{\DD}{{\cal D}}
\newcommand{\EE}{{\cal E}}
\newcommand{\FF}{{\cal F}}
\newcommand{\GG}{{\cal G}}
\newcommand{\HH}{{\cal H}}
\newcommand{\II}{{\cal I}}
\newcommand{\JJ}{{\cal J}}
\newcommand{\KK}{{\cal K}}
\newcommand{\LL}{{\cal L}}
\newcommand{\MM}{{\cal M}}
\newcommand{\NN}{{\cal N}}
\newcommand{\OO}{{\cal O}}
\newcommand{\PP}{{\cal P}}
\newcommand{\QQ}{{\cal Q}}
\newcommand{\QQQ}{{\cal Q}}
\newcommand{\RR}{{\cal R}}
\newcommand{\SSS}{{\cal S}}
\newcommand{\TT}{{\cal T}}
\newcommand{\TTT}{{\cal T}}
\newcommand{\UU}{{\cal U}}
\newcommand{\VV}{{\cal V}}
\newcommand{\WW}{{\cal W}}
\newcommand{\XX}{{\cal X}}
\newcommand{\YY}{{\cal Y}}
\newcommand{\ZZ}{{\cal Z}}
\newcommand{\SRB}{\mathscr{E}}
\newcommand{\Mix}{{\rm Mix}}

\newcommand{\fa}{{\mathfrak a}}
\newcommand{\fg}{{\mathfrak g}}
\newcommand{\fu}{{\mathfrak u}}
\newcommand{\fD}{{\mathfrak D}}
\newcommand{\fF}{{\mathfrak F}}
\newcommand{\fG}{{\mathfrak G}}
\newcommand{\fN}{{\mathfrak N}}
\newcommand{\fR}{{\mathfrak R}}
\newcommand{\fU}{{\mathfrak U}}
\newcommand{\fX}{{\mathfrak X}}


\newcommand{\Aa}{{\mathbb{A}}}
\newcommand{\Bb}{{\mathbb{B}}}
\newcommand{\Cc}{{\mathbb{C}}}
\newcommand{\Dd}{{\mathbb{D}}}
\newcommand{\Ee}{{\mathbb{E}}}
\newcommand{\Ff}{{\mathbb{F}}}
\newcommand{\Gg}{{\mathbb{G}}}
\newcommand{\Hh}{{\mathbb{H}}}
\newcommand{\Ii}{{\mathbb{I}}}
\newcommand{\Jj}{{\mathbb{J}}}
\newcommand{\Kk}{{\mathbb{K}}}
\newcommand{\Ll}{{\mathbb{L}}}
\newcommand{\Mm}{{\mathbb{M}}}
\newcommand{\Nn}{{\mathbb{N}}}
\newcommand{\Oo}{{\mathbb{O}}}
\newcommand{\Pp}{{\mathbb{P}}}
\newcommand{\Qq}{{\mathbb{Q}}}
\newcommand{\Rr}{{\mathbb{R}}}
\newcommand{\Ss}{{\mathbb{S}}}
\newcommand{\Tt}{{\mathbb{T}}}
\newcommand{\Uu}{{\mathbb{U}}}
\newcommand{\Vv}{{\mathbb{V}}}
\newcommand{\Ww}{{\mathbb{W}}}
\newcommand{\Xx}{{\mathbb{X}}}
\newcommand{\Yy}{{\mathbb{Y}}}
\newcommand{\Zz}{{\mathbb{Z}}}


\newcommand{\Q}{{\mathbb{Q}}}
\newcommand{\Z}{{\mathbb{Z}}}

\DeclareMathOperator{\sint}{int}
%%%%%%%%%%%%%%%%%%%%%%%%%%%%%%%%%%%%%%%%%%%%%

\def\H{{\mathfrak H}}
\def\Iso{\operatorname{Iso}}
\def\e{\mathrm{e}}
%\def\i{\mathrm{i}}
\def\cc{\operatorname{cc}}
\def\diag{\operatorname{diag}}
\def\dist{\operatorname{dist}}
\def\error{\operatorname{error}}
%\def\id{\operatorname{id}}
\def\j{\operatorname{j{}}}
\def\C{\operatorname{C{}}}
\def\G{\operatorname{G{}}}
\def\L{\operatorname{L{}}}
\def\M{\operatorname{M{}}}
\def\Mat{\operatorname{Mat}}
\def\GL{\operatorname{GL}}
\def\Op{\operatorname{Op}}
\def\sOp{\widetilde{\operatorname{Op}}}
\def\PGL{\operatorname{PGL}}
\def\Res{\operatorname{Res}}
\def\S{\operatorname{S{}}}
\def\Sp{\operatorname{Sp}}
\def\Sw{{\mathcal S}}
\def\SL{\operatorname{SL}}
\def\sl{\operatorname{sl}}
\def\SO{\operatorname{SO}}
\def\PSL{\operatorname{PSL}}
\def\O{\operatorname{O{}}}
\def\T{\operatorname{T{}}}
\def\tr{\operatorname{tr}}
\def\sgn{\operatorname{sgn}}
\def\supp{\operatorname{supp}}
\def\meas{\operatorname{meas}}
\def\Leb{\operatorname{Leb}}
\def\Var{\operatorname{Var}}
\def\Vol{\operatorname{Vol}}
\def\Area{\operatorname{Area}}
\def\ord{\operatorname{ord}}
\def\Prob{\operatorname{Prob}}
\def\Ad{\operatorname{Ad}}
\def\ad{\operatorname{ad}}

\def\Per{\operatorname{Per}}

\def\GamG{\Gamma\backslash G}
\def\SLZ{\SL(2,\Zz)}
\def\SLR{\SL_2(\Rr)}
\def\SLC{\SL_2(\Cc)}
\def\SOR{\SO(2)}

\def\Re{\operatorname{Re}}
\def\Im{\operatorname{Im}}

%%%%% Pedro's Macro
\newcommand{\Mod}[1]{\left\vert{#1}\right\vert}
\newcommand{\dem}{ \par\medbreak\noindent{\bf
Proof. }\enspace} 
\newcommand{\cqd}{\hfill
$\sqcup\!\!\!\!\sqcap\bigskip$}
\newcommand{\nrm}[1]{\left\|#1\right\|}


%\def\trans{\,^\top\!}

%%%%%%%%%%%%%%%%%%%%%%%%%%%%%%%%%%%%%%%%%%%%%%%%%%%%%%%%%%%%%%%%%%%%%

\newcommand{\abs}[1]{\left| #1\right|}
\newcommand{\norm}[1]{\left\| #1\right\|}
%\newcommand{\pde}[2]{\frac{\partial #1}{\partial #2}}

%\newcommand{\spec}[1]{\operatorname{spec}\left( #1\right)}
%\newcommand{\tr} {\operatorname{tr}}
%\newcommand{\diag} {\operatorname{diag}}
\newcommand{\trans} {\,^\top\!}
\newcommand{\conj} {\overline}
\newcommand{\para} {\parallel}

\newcommand{\id}  {\operatorname{id}}
\newcommand{\im}  {\operatorname{Im}}
\newcommand{\re}  {\operatorname{Re}}
%\newcommand{\up}  {\uparrow}
%\newcommand{\down}{\downarrow}
%\newcommand{\sgn} {\operatorname{sgn}}
%\newcommand{\const}{\operatorname{cst}}

\newcommand{\Int} {\operatorname{int}}
%\newcommand{\Ext} {\operatorname{Ext}}
%\newcommand{\Span}{\operatorname{span}}
%\newcommand{\perm}{\operatorname{Perm}}
%\newcommand{\std} {\operatorname{std}}
%\newcommand{\fix} {\operatorname{Fix}}
%\newcommand{\cl}  {\operatorname{ cl}}

\newcommand{\rot}{\operatorname{rot}}
\newcommand{\Homeo}{\operatorname{Homeo}} %{\mbox{\rm Homeo\,}}
\newcommand{\Diff} {\operatorname{Diff}}  %{\mbox{\rm Diff\,}}
\newcommand{\vf}   {\operatorname{Vect}} % {\XX}
%\newcommand{\diffo}{\text{Diff}_0\,}
%\newcommand{\vol}  {\operatorname{Diff}_{\text{vol}}\,}
%\newcommand{\Vol}  {\operatorname{vol}}
\newcommand{\Symp} {\operatorname{Symp}}
%\newcommand{\sympo}{\mbox{\rm Symp_0\,}} % {\text{Symp}_0\,}
\newcommand{\Ham}  {\operatorname{Ham}}
\newcommand{\Sym} {\operatorname{Sym}}
\newcommand{\Skew} {\operatorname{Skew}}
%\newcommand{\Sp}   {\operatorname{Sp}}
%\newcommand{\lie}  {\mbox{\rm Lie\,}}

%\newcommand{\GL}   {\operatorname{GL}}
%\newcommand{\SL}   {\operatorname{SL}}
%%\newcommand{\SU}   {\operatorname{SU}}
%\newcommand{\SO}   {\operatorname{SO}}

\newcommand{\DC}   {\operatorname{DC}}
\newcommand{\LC}   {\operatorname{LC}}

\newcommand{\te}[1]{\quad\text{#1}\quad}
%\newcommand{\comment}[1]{}

\newcommand{\Teich}{\mathscr{T}}
\newcommand{\TTeich}{{\rm T}\mathscr{T}}
\newcommand{\Bundle}{\mathscr{B}}
\newcommand{\Afrk}{\mathfrak{A}}
\newcommand{\Bfrk}{\mathfrak{B}}
\newcommand{\Cfrk}{\mathfrak{C}}
\newcommand{\Dfrk}{\mathfrak{D}}
\newcommand{\Ffrk}{\mathfrak{F}}
\newcommand{\Hfrk}{\mathfrak{H}}
\newcommand{\Lfrk}{\mathfrak{L}}
\newcommand{\Ofrk}{\mathfrak{O}}
\newcommand{\Qfrk}{\mathfrak{Q}}
\newcommand{\Rfrk}{\mathfrak{R}}
\newcommand{\Tfrk}{\mathfrak{T}}
\newcommand{\Ufrk}{\mathfrak{U}}
\newcommand{\Ascr}{\mathscr{A}}
\newcommand{\Bscr}{\mathscr{B}}
\newcommand{\Cscr}{\mathscr{C}}
\newcommand{\Escr}{\mathscr{E}}
\newcommand{\Zscr}{\mathscr{Z}}
\newcommand{\Gscr}{\mathscr{G}}
\newcommand{\Nscr}{\mathscr{N}}
\newcommand{\Mscr}{\mathscr{M}}
\newcommand{\Kscr}{\mathscr{K}}
\newcommand{\Iscr}{\mathscr{I}}
\newcommand{\Jscr}{\mathscr{J}}
\newcommand{\Dscr}{\mathscr{D}}
\newcommand{\Fscr}{\mathscr{F}}
\newcommand{\Hscr}{\mathscr{H}}
\newcommand{\Lscr}{\mathscr{L}}
\newcommand{\Oscr}{\mathscr{O}}
\newcommand{\Qscr}{\mathscr{Q}}
\newcommand{\Pscr}{\mathscr{P}}
\newcommand{\Rscr}{\mathscr{R}}
\newcommand{\Sscr}{\mathscr{S}}
\newcommand{\Tscr}{\mathscr{T}}
\newcommand{\Uscr}{\mathscr{U}}
\newcommand{\Vscr}{\mathscr{V}}
\newcommand{\Wscr}{\mathscr{W}}
\newcommand{\Xscr}{\mathscr{X}}
\newcommand{\Grp}{{\rm G}}
\newcommand{\Grass}{\mathscr{G}}
\newcommand{\Sperm}{{\rm S}}
\newcommand{\SJ}{{\rm S}^{sp}}
\newcommand{\Flag}{\Fscr}
\newcommand{\Fliso}{\Fscr^{sp}}
\newcommand{\Grptil}{\widetilde{\rm G}}
\newcommand{\Vfrk}{\mathfrak{V}}
\newcommand{\Xfrk}{\mathfrak{X}}
\newcommand{\Gfrk}{\mathfrak{G}}
\newcommand{\Pfrk}{\mathfrak{P}}
\newcommand{\Mfrk}{\mathfrak{M}}
\newcommand{\Ical}{\mathcal{I}}
\newcommand{\Grad}{{\rm Grad}}
\newcommand{\Sone}{[0,1]}

\newcommand{\emb}[1]{\BB^{#1}}
\newcommand{\ff}{\mathbb{II}}
\newcommand{\bepsilon}{\boldsymbol{\varepsilon}}
\newcommand{\vecomega}{{\boldsymbol{\omega}}}
\newcommand{\vecx}{{\boldsymbol{x}}}
\newcommand{\vecy}{{\boldsymbol{y}}}
\newcommand{\vecB}{{\boldsymbol{B}}}

\begin{document}
\title%[Positive Lyapunov Exponents versus Integrability]
%{On the positivity of Lyapunov exponents for random perturbations of billiards and standard maps}
[Lyapunov exponents and rigidity]{Lyapunov exponents and rigidity in random billiards and standard maps}

\date{\today}

\author[Del Magno]{Gianluigi Del Magno}
\address{Dipartimento di Matematica, Universit\`a di Pisa \\
Largo Bruno Pontecorvo 5, 56127 Pisa, Italy}
\email{gianluigi.delmagno@unipi.it}

\author[Lopes Dias]{Jo\~{a}o Lopes Dias}
\address{Universidade de Lisboa, ISEG Lisbon School of Economics \& Management, ISEG Research \\
Rua do Quelhas 6, 1200-781 Lisboa, Portugal}
\email{jldias@iseg.ulisboa.pt}

\author[Gaiv\~ao]{Jos\'e Pedro Gaiv\~ao}
\address{Universidade de Lisboa, ISEG Lisbon School of Economics \& Management, ISEG Research \\
Rua do Quelhas 6, 1200-781 Lisboa, Portugal}
\email{jpgaivao@iseg.ulisboa.pt}


\begin{abstract}
We study random dynamical systems generated by measure-preserving maps, possibly with singularities. For this class of systems, we establish an invariance principle: if all Lyapunov exponents coincide almost everywhere, then there exists an invariant measurable family of probability measures on the projective tangent bundle for the projective cocycle induced by the derivative. We apply this principle to random additive perturbations of two classes of maps: billiard maps for strictly convex tables on surfaces of constant curvature and generalized standard maps. In both settings, we obtain rigidity results: the Lyapunov exponents vanish almost everywhere precisely for billiards in geodesic disks and generalized standard maps with constant potential; in all other cases, they are nonzero almost everywhere.


% We study random dynamical systems generated by maps with singularities that preserve a common probability measure.
% We establish an invariance principle stating that, if the maximal and minimal Lyapunov exponents coincide almost everywhere, then there exists a probability measure on the projective tangent bundle, projecting to the common invariant measure, that is invariant under the projectivized derivative of almost every map generating the random system. For random additive perturbations of volume-preserving toral maps with absolutely continuous noise, vanishing Lyapunov exponents impose strong geometric and algebraic restrictions on the derivatives of the underlying deterministic map. We apply these results to two classes of conservative systems. For random additive perturbations of the billiard map associated with a strictly convex table on a surface of constant curvature, we prove that the Lyapunov exponents are non-zero almost everywhere unless the table is a geodesic disk, in which case they vanish almost everywhere. For random additive perturbations of generalized standard maps, we prove that the Lyapunov exponents are non-zero almost everywhere unless the potential is constant, in which case they vanish almost everywhere.
\end{abstract}

\maketitle
\tableofcontents

%%%%%%%%%%%%%%%%%%%%%%%%%%%%%%%%%%%%%%
\section{Introduction}
\label{sec:introduction}
Numerical simulations indicate that many volume-preserving maps exhibit a mixed phase space, with coexisting regions of regular and chaotic dynamics; see~\cite{Strelcyn1991} and the references therein. This behavior occurs in a broad class of systems, including area-preserving twist maps and billiards in convex tables, and in many examples suggests the presence of positive Lyapunov exponents on a set of positive volume. Despite substantial numerical and heuristic evidence, establishing this property rigorously for deterministic conservative systems remains a difficult problem.

This difficulty motivates the study of conservative systems in the presence of noise. Random perturbations are expected to weaken invariant structures and improve ergodic properties, thereby promoting exponential instability. In this paper, we study random additive perturbations of conservative maps and ask whether they have non-zero Lyapunov exponents almost everywhere. When the exponents vanish, we ask whether the underlying deterministic system has special geometric or algebraic properties. This is the rigidity problem at the center of the paper. 
%For the two classes considered here, the answer is affirmative: vanishing Lyapunov exponents characterize the exceptional integrable members of the corresponding families.

A key tool in our analysis is an invariance principle for random dynamical systems generated by maps with singularities. This terminology refers to maps that are $C^1$ diffeomorphisms on an open set of full measure, as in the framework of~\cite{Katok:1986tg}.

Let $M$ be a Riemannian manifold, let $\nu$ be a probability measure on a space $X$, and let
\[
f_x\colon M\to M,
\qquad x\in X,
\]
be a measurable family of maps preserving a common probability measure $m$ on $M$. Set
\[
\Omega=X^{\Nn},
\qquad
\rho_\nu=\nu^{\Nn},
\]
write $\omega=(\omega_0,\omega_1,\ldots) \in \Omega$, and denote by $\sigma\colon\Omega\to\Omega$ the left shift. The associated random dynamical system is the skew product
\[
F\colon\Omega\times M\to\Omega\times M,
\qquad
F(\omega,y)
=
\bigl(\sigma(\omega),f_{\omega_0}(y)\bigr).
\]
The random compositions generated by the family $(f_x)_{x\in X}$ are
\[
F_\omega^0=\operatorname{id}_M,
\qquad
F_\omega^n
=
f_{\omega_{n-1}}\circ\cdots\circ f_{\omega_0},
\qquad n\geq1.
\]

We assume that each $f_x$ restricts to a $C^1$ diffeomorphism on an open set of full $m$-measure and that the derivative cocycle satisfies suitable integrability conditions. We denote by $\lambda_F^+$ and $\lambda_F^-$ the maximal and minimal Lyapunov exponents of $F$, respectively. The precise assumptions and definitions are given in Section~\ref{sec:random maps}. Under these hypotheses, we prove the following result.

\begin{theorem}
Suppose that $\lambda_F^-(\omega,y) = \lambda_F^+(\omega,y)$ for $(\rho_\nu\times m)$-almost every $(\omega,y)\in\Omega\times M$. Then there exists a probability measure $\eta$ on the projective tangent bundle $PM$ such that $\pi_*\eta=m$, where $\pi\colon PM\to M$ is the bundle projection, and
\[
(Df_x)_*\eta=\eta
\qquad \text{for } \nu \text{-a.e. } x\in X.
\]
Equivalently, if $\{\eta_y:y\in M\}$ denotes the disintegration of $\eta$ over $m$, then
\[
Df_x(y)_*\eta_y
=
\eta_{f_x(y)}
\]
for $(\nu\times m)$-almost every $(x,y)\in X\times M$.
\end{theorem}

Following Avila and Viana~\cite{Avila:2010aa}, we refer to this result as an invariance principle. It provides the basic rigidity mechanism used throughout the paper: coincidence of the extremal Lyapunov exponents forces the existence of a non-random invariant structure on the projective tangent bundle. The result is obtained by applying a theorem of Ledrappier~\cite{Ledrappier:1986ue} for linear cocycles over stationary Markov chains to the derivative cocycle of the random system. The underlying principle goes back to Furstenberg's theory of random matrix products~\cite{Furstenberg:1963uq}. Analogous invariance principles for random diffeomorphisms were proved by Carverhill~\cite{carverhill1987furstenberg} and Baxendale~\cite{Baxendale:1989aa}.

We next specialize the invariance principle to random additive perturbations of conservative maps on the torus. Let
\[
g\colon\Tt^d\to\Tt^d
\]
be a map that preserves the normalized Haar measure $m$ and satisfies the regularity and integrability assumptions introduced in Definition~\ref{def:toral_map}. In this setting, $X=\Tt^d$ and $\Omega=(\Tt^d)^{\Nn}$. For each $x\in\Tt^d$, define
\[
f_x(y)=g(y)+x.
\]
Given an i.i.d.\ sequence $\omega=(\omega_n)_{n\geq0} \in \Omega$ of $\Tt^d$-valued random variables with common distribution $\nu$, consider the corresponding random compositions $F_\omega^n$. 

For these systems, the sum of the Lyapunov exponents vanishes almost everywhere. Consequently, $\lambda_F^-=\lambda_F^+$ almost everywhere if and only if $\lambda_F^+=0$ almost everywhere. The invariance principle therefore has the following consequence. We write $\SL^{\pm}(d,\Rr)$ for the set of real $d \times d$ matrices with determinant $\pm 1$.

\begin{theorem}
Let $F$ be a $\nu$-random additive perturbation of a toral map
$g\colon\Tt^d\to\Tt^d$, and let $N\subset\Tt^d$ be an open set of full
$m$-measure such that $g|_N\colon N\to g(N)$ is a $C^1$ diffeomorphism. Suppose that $\lambda_F^+(\omega,y)=0$ for $(\rho_\nu\times m)$-almost every
$(\omega,y)\in\Omega\times\Tt^d$. Then there exists a measurable family
of probability measures $\{\eta_y:y\in\Tt^d\}$ on $\mathbb P^{d-1}$ such
that
\[
Dg(y)_*\eta_y=\eta_{g(y)+x}
\]
for $m$-almost every $y\in\Tt^d$ and $\nu$-almost every $x\in\Tt^d$. If, in addition, $\nu$ is absolutely continuous with respect to $m$, then
there exists a probability measure $\eta$ on $\mathbb P^{d-1}$ such that
\[
h_*\eta=\eta
\]
for every $h\in\operatorname{supp}\bigl((Dg)_*(m|_N)\bigr)$. Hence, the support of $(Dg)_*(m|_N)$ is either contained in a
compact subgroup of $\operatorname{SL}^{\pm}(d,\Rr)$ or preserves a finite family of proper subspaces of $\Rr^d$.
\end{theorem}

Thus, for random additive perturbations with absolutely continuous noise, vanishing Lyapunov exponents force the derivatives of the underlying deterministic map to preserve a common projective probability measure. 

We apply this rigidity mechanism to two fundamental classes of conservative dynamical systems: billiards in convex tables and generalized standard maps. For random additive perturbations of convex billiards, we prove that the Lyapunov exponents vanish almost everywhere precisely if and only if the table is a geodesic disk. For random additive perturbations of generalized standard maps, they vanish almost everywhere if and only if the potential is constant.

We begin with billiards on surfaces of constant curvature. Let $S$ be one of the standard Riemannian surfaces of constant curvature: the Euclidean plane, the sphere, or the hyperbolic plane. Let $D\subset S$ be a strictly convex domain whose boundary $\partial D$ is a $C^2$ simple closed curve of positive geodesic curvature and unit length. We refer to such a domain $D$ as a convex table. A billiard in $D$ describes the motion of a point particle travelling along geodesics of $S$ inside $D$ and undergoing elastic reflections at $\partial D$.

The dynamics is encoded by the billiard map $\Phi\colon V\to V$, a twist map on the collision space $V=S^1\times[-1,1]$. Each point $(s,r)\in V$ represents a collision with the boundary: $s$ is the arc-length parameter of the collision point on $\partial D$, and $r=-\cos\theta$, where $\theta\in[0,\pi]$ is the reflection angle measured with respect to the tangent to $\partial D$ at $s$. The map $\Phi$ sends each collision to the next collision along the trajectory. After the identification described in Section~\ref{sec:billiards}, the billiard map induces a map $T\colon\Tt^2\to\Tt^2$ preserving the normalized Haar measure $m$.

Despite extensive study, it remains an open problem whether billiards in strictly convex tables with smooth boundary can exhibit positive metric entropy with respect to $m$. All known examples of billiards with this property involve tables that are nor strictly convex nor smooth; see, for instance, \cite{Bunimovich:79,Wojtkowski:86,Donnay:91,Markarian:94,Gutkin:99}. Although generic convex billiards have positive topological entropy \cite{BessaDelMagnoLopesDiasGaivaoTorres24,Cheng2004400,ZHANG2017793}, this does not by itself provide examples with positive metric entropy with respect to $m$.

The following theorem establishes a sharp rigidity dichotomy for random additive perturbations of convex billiards, distinguishing geodesic disks from all other convex tables. Here $T$ denotes the billiard map on $\Tt^2$, the maps
\[
f_x(y)=T(y)+x, \qquad x \in \Tt^2,
\]
are its additive perturbations, and $F_\omega^n$ denotes the corresponding random composition. Additive perturbations of the billiard map admit a natural geometric interpretation. Given $x=(\bar s,\bar r)\in\Tt^2$ and an initial collision $(s_0,r_0)$, one first computes the next collision $(s_1,r_1)=T(s_0,r_0)$
and then translates it by $x$ in phase space, shifting the arc-length coordinate by $\bar s$ and the $r$-coordinate by $\bar r$.

\begin{theorem}
\label{thm:main}
Let $D$ be a convex billiard table on a surface $S$ of constant curvature. Assume that the common distribution $\nu$ of the random variables $(\omega_n)_{n\geq0}$ is absolutely continuous with respect to $m$. Then the Lyapunov exponents of the random compositions $F_\omega^n$ are non-zero almost everywhere if $D$ is not a geodesic disk, and vanish almost everywhere otherwise.
\end{theorem}

This theorem reflects the exceptional symmetry of circular billiards on surfaces of constant curvature. If $D$ is a geodesic disk, then the phase space is foliated by non-contractible invariant circles \cite{Bialy:1993ty,dosSantos:2017}, a property commonly referred to as total integrability. Under an additive perturbation, this foliation remains globally invariant: individual circles need not be invariant, but each circle is mapped onto another circle of the same foliation. Consequently, the derivative cocycle has no exponential growth, and the Lyapunov exponents vanish almost everywhere.

The converse is the rigidity direction of the theorem. Suppose that the Lyapunov exponents vanish. The invariance principle and its consequences for additive perturbations yield a common invariant projective probability measure. The factorization of the derivative cocycle, together with its behavior near tangential collisions, forces this measure to be the Dirac mass at the horizontal direction. Hence the horizontal direction is invariant under the derivative of the billiard map. This implies that the billiard foliation by horizontal circles is invariant and ultimately forces the boundary of the table to have constant geodesic curvature. Therefore $D$ is a geodesic disk.

Theorem~\ref{thm:main} considerably broadens the scope of the result of Zhang and Nguyen~\cite{NguyenZhang2026}, who proved the existence of non-zero Lyapunov exponents for random additive perturbations of billiards in non-circular elliptical tables and lemon-shaped tables in the Euclidean plane. It also provides a new rigidity characterization of circular billiards among convex billiards on surfaces of constant curvature (cf.~\cite{Bialy:1993ty} and~\cite{CzudekDeSimoiGadPoon26}). 

The second class of conservative systems considered in this paper consists of generalized standard maps. Given a $C^1$ function
\[
V\colon\Tt^1\to\Rr,
\]
the associated generalized standard map $g_V\colon\Tt^2\to\Tt^2$ is
\[
g_V(y_1,y_2)
=
\bigl(y_1+y_2+V(y_1),\,y_2+V(y_1)\bigr)
\pmod 1.
\]
For $V(y_1)=K\sin(2\pi y_1)$, this is the standard map introduced by Chirikov~\cite{CHIRIKOV1979263}; the case $K=0$ is completely integrable.

Generalized standard maps arise naturally in Hamiltonian dynamics as models of periodically kicked systems and play a central role in the study of the transition from integrability to chaos. Nevertheless, proving the existence of non-zero Lyapunov exponents for non-integrable standard maps and their generalizations remains notoriously difficult.

The following theorem is the counterpart of Theorem~\ref{thm:main} for generalized standard maps. For each $x\in\Tt^2$, define
\[
f_x(y)=g_V(y)+x,
\]
and let $F_\omega^n$ denote the corresponding random composition. The theorem establishes the corresponding sharp rigidity dichotomy, distinguishing constant potentials from non-constant ones. In this setting, the conclusion holds both for absolutely continuous noise and for a class of degenerate noise distributions.

We call a Borel probability measure $\nu$ on $\Tt^2$ a \emph{vertical line measure} if
\[
\nu=\delta_0\times\chi,
\]
where $\delta_0$ is the Dirac measure at $0\in\Tt^1$ and
$\chi=h\,m_1$ is absolutely continuous with respect to the normalized Haar measure $m_1$ on $\Tt^1$.

\begin{theorem}
\label{thm:standard}
Let $g_V\colon\Tt^2\to\Tt^2$ be a generalized standard map. Assume that the common distribution $\nu$ of the random variables $(\omega_n)_{n\geq0}$ is either absolutely continuous with respect to $m$ or a vertical line measure. Then the Lyapunov exponents of the random compositions $F_\omega^n$ are non-zero almost everywhere if $V$ is non-constant, and vanish almost everywhere otherwise.
\end{theorem}

The derivative matrices of a non-constant generalized standard map contain both hyperbolic and elliptic elements, which excludes the alternatives compatible with an invariant projective probability measure. For vertical line noise, the random variables are grouped in pairs, reducing the problem to an additive perturbation with an absolutely continuous noise distribution.

When $\nu$ is absolutely continuous, the result follows %by combining ergodicity with the projective rigidity supplied by 
from the invariance principle. If $V$ is non-constant, among the derivative matrices, there are both hyperbolic and elliptic elements, which excludes the alternatives compatible with an invariant projective probability measure. For vertical line noise, the random variables are grouped in pairs, reducing the problem to a random additive perturbation of $g_V^2$ with an absolutely continuous noise distribution.

For random standard maps, Blumenthal, Xue, and Young~\cite{BlumenthalXueYoung} obtained quantitative lower bounds for the maximal Lyapunov exponent under sufficiently strong random additive perturbations. In contrast, the present result does not require the noise to be large and gives a sharp characterization of the vanishing of the Lyapunov exponents.

The proofs of the two theorems follow a common rigidity scheme. First, we prove that the stationary measure is ergodic for the relevant random additive perturbations. This yields a dichotomy for the maximal Lyapunov exponent: it either vanishes almost everywhere or is positive almost everywhere. In the vanishing case, the invariance principle produces a measurable invariant family of projective measures. The additive noise promotes this family to a common projective measure, and the structure of the derivative cocycle turns that measure into an invariant geometric structure for the deterministic map. The final rigidity step classifies the maps carrying this structure: the billiard table must be a geodesic disk, and the potential of the generalized standard map must be constant.

The paper is organized as follows. Section~\ref{sec:random maps} establishes a general invariance principle for random piecewise smooth maps and derives several consequences needed later in the paper. Section~\ref{se:ergodicity} proves the unique ergodicity of random additive perturbations and the equidistribution of their random orbits in two general settings. These settings cover the random additive perturbations of billiard maps and generalized standard maps treated in the final sections. Section~\ref{sec:billiards} provides background on convex billiards on surfaces of constant curvature. The proof of Theorem~\ref{thm:main} is given in Section~\ref{ss:random_billiards}. Finally, Section~\ref{sec:random_standard} is devoted to random generalized standard maps and contains the proof of
Theorem~\ref{thm:standard}.


%--- Old version
%Numerical simulations indicate that many volume-preserving maps exhibit a mixed phase space, with coexisting regions of regular and chaotic dynamics; see~\cite{Strelcyn1991} and the references therein. This behavior occurs in a broad class of systems, including area-preserving twist maps and billiards in convex tables, and in many examples suggests the presence of positive Lyapunov exponents on a set of positive volume. Despite substantial numerical and heuristic evidence, establishing this property rigorously for deterministic conservative systems remains a difficult problem.

%This difficulty motivates the study of conservative systems in the presence of noise. Random perturbations are expected to weaken invariant structures and improve ergodic properties, thereby promoting exponential instability. In this paper, we study random additive perturbations of conservative maps and ask whether they have non-zero Lyapunov exponents almost everywhere.  when they do not, whether the persistence of vanishing Lyapunov exponents reflects integrability of the underlying deterministic system.

%A key tool in our analysis is an invariance principle for random dynamical systems generated by maps with singularities. The terminology refers to maps that are $C^1$ diffeomorphisms on open set of full measure, as in the framework of~\cite{Katok:1986tg}.

%Let $M$ be a Riemannian manifold, let $\nu$ be a probability measure on a space $X$, and let
%\[
%f_x\colon M\to M,
%\qquad x\in X,
%\]
%be a measurable family of maps preserving a common probability measure $m$ on $M$. Set
%\[
%\Omega=X^{\Nn},
%\qquad
%\rho_\nu=\nu^{\Nn},
%\]
%write $\omega=(\omega_0,\omega_1,\ldots) \in \Omega$, and denote by $\sigma\colon\Omega\to\Omega$ the left shift. The associated random dynamical system is the skew product
%\[
%F\colon\Omega\times M\to\Omega\times M,
%\qquad
%F(\omega,y)
%=
%\bigl(\sigma(\omega),f_{\omega_0}(y)\bigr).
%\]
%The random compositions generated by the family $(f_x)_{x\in X}$ are
%\[
%F_\omega^0=\operatorname{id}_M,
%\qquad
%F_\omega^n
%=
%f_{\omega_{n-1}}\circ\cdots\circ f_{\omega_0},
%\qquad n\geq1.
%\]

%We assume that each $f_x$ restricts to a $C^1$ diffeomorphism on an open set of full $m$-measure and that the derivative cocycle satisfies suitable integrability conditions. We denote by $\lambda_F^+$ and $\lambda_F^-$ the maximal and minimal Lyapunov exponents of $F$, respectively. The precise assumptions and definitions are given in Section~\ref{sec:random maps}. Under these hypotheses, we prove the following result.

%\begin{theorem}
%Suppose that
%\[
%\lambda_F^-(\omega,y)
%=
%\lambda_F^+(\omega,y)
%\]
%for $(\rho_\nu\times m)$-almost every $(\omega,y)\in\Omega\times M$. Then there exists a probability measure $\eta$ on the projective tangent bundle $PM$ such that $\pi_*\eta=m$, where $
%\pi\colon PM\to M$ is the bundle projection, and
%\[
%(Df_x)_*\eta=\eta
%\qquad \text{for } \nu \text{-a.e. } x\in X.
%\]
%Equivalently, if $\{\eta_y:y\in M\}$ denotes the disintegration of $\eta$ over $m$, then
%\[
%Df_x(y)_*\eta_y
%=
%\eta_{f_x(y)}
%\]
%for $(\nu\times m)$-almost every $(x,y)\in X\times M$.
%\end{theorem}

%Following Avila and Viana~\cite{Avila:2010aa}, we refer to this result as an invariance principle. It is obtained by applying a theorem of Ledrappier~\cite{Ledrappier:1986ue} for linear cocycles over stationary Markov chains to the derivative cocycle of the random system. The principle that coincidence of the extremal Lyapunov exponents yields an invariant probability measure on projective space goes back to Furstenberg's theory of random matrix products~\cite{Furstenberg:1963uq}. Analogous invariance principles for random diffeomorphisms were proved by Carverhill~\cite{carverhill1987furstenberg} and Baxendale~\cite{Baxendale:1989aa}.

%We next specialize the invariance principle to random additive perturbations of conservative maps on the torus. Let
%\[
%g\colon\Tt^d\to\Tt^d
%\]
%be a map preserving the normalized Haar measure $m$ and satisfying the regularity and integrability assumptions introduced in Definition~\ref{def:toral_map}. In this setting, $X=\Tt^d$ and $\Omega=(\Tt^d)^{\Nn}$. For each $x\in\Tt^d$, define
%\[
%f_x(y)=g(y)+x.
%\]
%Given an i.i.d.\ sequence $\omega=(\omega_n)_{n\geq0}$ of $\Tt^d$-valued random variables with common distribution $\nu$, consider the corresponding random compositions $F_\omega^n$. Since every $f_x$ preserves $m$, the measure $m$ is stationary for the associated Markov chain.

%For these systems, the sum of the Lyapunov exponents vanishes almost everywhere. Consequently, $\lambda_F^-=\lambda_F^+$ almost everywhere if and only if $
%\lambda_F^+=0$ almost everywhere. The invariance principle therefore has the following consequence. We write $\SL^{\pm}(d,\Rr)$ for the set of real $d \times d$ matrices with determinant $\pm 1$.

%\begin{theorem}
%Let $F$ be a $\nu$-random additive perturbation of a toral map
%$g\colon\Tt^d\to\Tt^d$, and let $N\subset\Tt^d$ be an open set of full
%$m$-measure such that $g|_N\colon N\to g(N)$ is a $C^1$ diffeomorphism. Suppose that
%\[
%\lambda_F^+(\omega,y)=0
%\]
%for $(\rho_\nu\times m)$-almost every
%$(\omega,y)\in\Omega\times\Tt^d$. Then there exists a measurable family
%of probability measures $\{\eta_y:y\in\Tt^d\}$ on $\mathbb P^{d-1}$ such
%that
%\[
%Dg(y)_*\eta_y=\eta_{g(y)+x}
%\]
%for $m$-almost every $y\in\Tt^d$ and $\nu$-almost every $x\in\Tt^d$. If, in addition, $\nu$ is absolutely continuous with respect to $m$, then
%there exists a probability measure $\eta$ on $\mathbb P^{d-1}$ such that
%\[
%h_*\eta=\eta
%\]
%for every $h\in\operatorname{supp}\bigl((Dg)_*(m|_N)\bigr)$. Hence, the support of $(Dg)_*(m|_N)$ is either contained in a
%compact subgroup of $\operatorname{SL}^{\pm}(d,\Rr)$ or preserves a finite family of proper subspaces of $\Rr^d$.
%\end{theorem}

%Thus, under absolutely continuous noise, for random additive perturbations, vanishing Lyapunov exponents force the derivative of the underlying deterministic map to preserve a common projective probability measure. This converts a dynamical condition into a geometric and algebraic property. It is this consequence of the invariance principle that is used in the applications below.

%We apply this principle to two fundamental classes of conservative dynamical systems: billiards in convex tables and generalized standard maps. For random additive perturbations of convex billiards, we prove that the Lyapunov exponents vanish almost everywhere precisely when the table is a geodesic disk. For random additive perturbations of generalized standard maps, they vanish almost everywhere when the potential is constant.

%We begin with billiards on surfaces of constant curvature. Let $S$ be one of the standard Riemannian surfaces of constant curvature: the Euclidean plane, the sphere, or the hyperbolic plane. Let $D\subset S$ be a strictly convex domain whose boundary $\partial D$ is a $C^2$ simple closed curve of positive geodesic curvature and unit length. We refer to such a domain $D$ as a convex table. A billiard in $D$ describes the motion of a point particle travelling along geodesics of $S$ inside $D$ and undergoing elastic reflections at $\partial D$.

%The dynamics is encoded by the billiard map $\Phi\colon V\to V$, a twist map on the collision space $V=S^1\times[-1,1]$. Each point $(s,r)\in V$ represents a collision with the boundary: $s$ is the arc-length parameter of the collision point on $\partial D$, and $r=-\cos\theta$, where $\theta\in[0,\pi]$ is the reflection angle measured with respect to the tangent to $\partial D$ at $s$. The map $\Phi$ sends each collision to the next collision along the trajectory. After the identification described in Section~\ref{sec:billiards}, the billiard map induces a map $T\colon\Tt^2\to\Tt^2$ preserving the normalized Haar measure $m$.

%Despite extensive study, it remains an open problem whether billiards in strictly convex tables with smooth boundary can exhibit positive metric entropy with respect to $m$. All known examples of billiards with this property involve tables that are neither strictly convex nor smooth; see, for instance, \cite{Bunimovich:79,Wojtkowski:86,Donnay:91,Markarian:94,Gutkin:99}. Although generic convex billiards have positive topological entropy \cite{BessaDelMagnoLopesDiasGaivaoTorres24,Cheng2004400,ZHANG2017793}, this does not by itself provide examples with positive metric entropy with respect to $m$.

%The following theorem establishes a sharp dichotomy for the Lyapunov exponents of random additive perturbations of convex billiards, distinguishing geodesic disks from all other convex tables. Here $T$ denotes the billiard map on $\Tt^2$, the maps
%\[
%f_x(y)=T(y)+x, \qquad x \in \Tt^2,
%\]
%are its additive perturbations, and $F_\omega^n$ denotes the corresponding random composition. Additive perturbations of the billiard map admit a natural geometric interpretation. Given $x=(\bar s,\bar r)\in\Tt^2$ and an initial collision $(s_0,r_0)$, one first computes the next collision $(s_1,r_1)=T(s_0,r_0)$
%and then translates it by $x$ in phase space, shifting the arc-length coordinate by $\bar s$ and the $r$-coordinate by $\bar r$.

%\begin{theorem}
%\label{thm:main}
%Let $D$ be a convex billiard table on a surface $S$ of constant curvature. Assume that the common distribution $\nu$ of the random variables $(\omega_n)_{n\geq0}$ is absolutely continuous with respect to $m$. Then the Lyapunov exponents of the random compositions $F_\omega^n$ are non-zero almost everywhere if $D$ is not a geodesic disk, and vanish almost everywhere otherwise.
%\end{theorem}

%This theorem reflects the exceptional symmetry of circular billiards on surfaces of constant curvature. If $D$ is a geodesic disk, then the phase space is foliated by non-contractible invariant circles \cite{Bialy:1993ty,dosSantos:2017}, a property commonly referred to as total integrability. Under an additive perturbation, this foliation remains globally invariant: individual circles need not be invariant, but each circle is mapped onto another circle of the same foliation. Consequently, the derivative cocycle has no exponential growth, and the Lyapunov exponents vanish almost everywhere.

%Conversely, suppose that the Lyapunov exponents vanish. The invariance principle and its consequences for additive perturbations yield an invariant projective probability measure. The singular behavior of the derivative near tangential collisions forces this measure to be concentrated on the horizontal direction. Hence the horizontal direction is invariant under the derivative of the billiard map. This implies that the billiard foliation by horizontal circles is invariant and ultimately forces the boundary of the table to have constant geodesic curvature. Therefore $D$ is a geodesic disk.

%Theorem~\ref{thm:main} considerably broadens the scope of the result of Zhang and Nguyen~\cite{NguyenZhang2026}, who proved the existence of non-zero Lyapunov exponents for random additive perturbations of billiards in non-circular elliptical tables and lemon-shaped tables in the Euclidean plane.
%It also provides another characterization of circular billiards among convex billiards on surfaces of constant curvature (cf.~\cite{Bialy:1993ty} and~\cite{CzudekDeSimoiGadPoon26}). 

%The second class of conservative systems considered in this paper consists of generalized standard maps. Given a $C^1$ function
%\[
%V\colon\Tt^1\to\Rr,
%\]
%the associated generalized standard map $g_V\colon\Tt^2\to\Tt^2$ is
%\[
%g_V(y_1,y_2)
%=
%\bigl(y_1+y_2+V(y_1),\,y_2+V(y_1)\bigr)
%\pmod 1.
%\]
%For $V(y_1)=K\sin(2\pi y_1)$, this is the standard map introduced by Chirikov~\cite{CHIRIKOV1979263}; the case $K=0$ is completely integrable.

%Generalized standard maps arise naturally in Hamiltonian dynamics as models of periodically kicked systems and play a central role in the study of the transition from integrability to chaos. Nevertheless, proving the existence of non-zero Lyapunov exponents for non-integrable standard maps and their generalizations remains notoriously difficult.

%The following theorem is the counterpart of Theorem~\ref{thm:main} for generalized standard maps. For each $x\in\Tt^2$, define
%\[
%f_x(y)=g_V(y)+x,
%\]
%and $F_\omega^n$ denotes the corresponding random composition. The theorem establishes a sharp dichotomy for the Lyapunov exponents of these random compositions, distinguishing constant potentials from non-constant ones. In this setting, the conclusion holds both for absolutely continuous noise and for a class of degenerate noise distributions.

%We call a Borel probability measure $\nu$ on $\Tt^2$ a \emph{vertical line measure} if
%\[
%\nu=\delta_0\times\chi,
%\]
%where $\delta_0$ is the Dirac measure at $0\in\Tt^1$ and
%$\chi=h\,m_1$ is absolutely continuous with respect to the normalized Haar measure $m_1$ on $\Tt^1$.

%\begin{theorem}
%\label{thm:standard}
%Let $g_V\colon\Tt^2\to\Tt^2$ be a generalized standard map. Assume that the common distribution $\nu$ of the random variables $(\omega_n)_{n\geq0}$ is either absolutely continuous with respect to $m$ or a vertical line measure. Then the Lyapunov exponents of the random compositions $F_\omega^n$ are non-zero almost everywhere if $V$ is non-constant, and vanish almost everywhere otherwise.
%\end{theorem}

%When $\nu$ is absolutely continuous, the result follows by combining ergodicity with the projective rigidity supplied by the invariance principle. The derivative matrices of a non-constant generalized standard map contain both hyperbolic and elliptic elements, which excludes the alternatives compatible with an invariant projective probability measure. For vertical line noise, the random variables are grouped in pairs, reducing the problem to an additive perturbation with an absolutely continuous noise distribution.

%For random standard maps, Blumenthal, Xue, and Young~\cite{BlumenthalXueYoung} obtained quantitative lower bounds for the maximal Lyapunov exponent under sufficiently strong random additive perturbations. In contrast, the present result does not require the noise to be large and gives a sharp characterization of the vanishing of the Lyapunov exponents.

%The proofs of the two theorems follow a common strategy. First, we prove that the stationary measure is ergodic for the relevant random additive perturbations. This yields a dichotomy for the maximal Lyapunov exponent: it either vanishes almost everywhere or is positive almost everywhere. Assuming that it vanishes, the invariance principle produces an invariant probability measure on projective space. The structure of the derivative cocycle then forces an invariant geometric structure for the underlying deterministic map. For billiards, this structure implies that the table is a geodesic disk; for generalized standard maps, it implies that the potential is constant.

%The paper is organized as follows. Section~\ref{sec:random maps} establishes a general invariance principle for random piecewise smooth maps and derives several consequences needed later in the paper. Section~\ref{se:ergodicity} proves the unique ergodicity of random additive perturbations and the equidistribution of their random orbits in two general settings. These settings cover the random additive perturbations of billiard maps and generalized standard maps treated in the final sections. Section~\ref{sec:billiards} provides background on convex billiards on surfaces of constant curvature. The proof of Theorem~\ref{thm:main} is given in Section~\ref{ss:random_billiards}. Finally, Section~\ref{sec:random_standard} is devoted to random generalized standard maps and contains the proof of
%Theorem~\ref{thm:standard}.

%\added{---Old}

%Numerical simulations indicate that generic volume-preserving maps often exhibit a mixed phase space, with coexisting regions of regular and chaotic dynamics (see~\cite{Strelcyn1991} and references therein). This behavior is observed in a broad class of systems, including area-preserving twist maps and billiards in convex tables, and in many examples suggests the presence of positive metric entropy with respect to Lebesgue measure. Despite substantial numerical and heuristic evidence, no general method is currently available to establish this property rigorously.

%\comment{New part here} In this paper, we study random additive perturbations of area-preserving maps on the 2-torus $\Tt^2$ and investigate whether randomness leads to non-zero Lyapunov exponents almost everywhere. More precisely, let $f \colon \Tt^2 \to \Tt^2$ be a map preserving the normalized Lebesgue measure $m$ on $\Tt^2$. For each $x \in \Tt^2$, define the additive perturbation of $f$ by
%\[
%f_x(y)=f(y)+x \pmod 1.
%\]
%Let $(\omega_n)_{n\ge 0}$ be an i.i.d.\ sequence of $\Tt^2$-valued random variables with common distribution $\nu$, and consider the random dynamical system generated by the maps
%\[
%F_\omega^0=\operatorname{id},
%\qquad
%F_\omega^n = f_{\omega_{n-1}}\circ\cdots\circ %f_{\omega_0}.
%\]
%Since each map $f_x$ preserves $m$, the measure $m$ is stationary for the Markov chain on $\Tt^2$ generated by the random iterates $F_\omega^n$.

%Additive randomness is expected to weaken invariant structures and improve ergodic properties, thereby promoting exponential instability and positive metric entropy. This leads to the following question: do conservative random perturbations typically generate non-zero Lyapunov exponents almost everywhere, and if not, does the persistence of vanishing Lyapunov exponents characterize integrability of the underlying deterministic system?

%We address this question for two fundamental classes of conservative dynamical systems: billiards in convex tables and standard maps.

%We begin with billiards on surfaces of constant curvature. Let $S$ be a Riemannian surface of constant curvature, and let $D\subset S$ be a strictly convex domain whose boundary $\partial D$ is a $C^2$ simple closed curve with positive geodesic curvature. We refer to such a domain $D$ as a \emph{convex table}. A billiard in $D$ describes the motion of a point particle traveling along geodesics of $S$ inside $D$ and undergoing elastic reflections at $\partial D$.

%The dynamics is encoded by the \emph{billiard map} $\Phi \colon V\to V$, a twist map on the collision space $V=S^1\times[-1,1]$, where $S^1=\Rr/(L\Zz)$ and $L$ denotes the length of $\partial D$. Each point $(s,r)\in V$ represents a collision with the boundary: $s$ is the arc-length parameter of the collision point on $\partial D$, and $r=-\cos\theta$, where $\theta\in[0,\pi]$ is the reflection angle measured with respect to the tangent to $\partial D$ at $s$. The map $\Phi$ sends each collision $(s,r)$ to the next collision along the trajectory. After the identification described in Section~\ref{sec:billiards}, the billiard map induces a map $f\colon\Tt^2\to\Tt^2$ preserving $m$.

%Despite extensive study, it remains an open problem whether billiards in strictly convex tables with smooth boundary can exhibit positive metric entropy with respect to $m$. All known examples of billiards with this property involve tables that are either not strictly convex or not smooth; see, for instance, \cite{Bunimovich:79,Wojtkowski:86,Donnay:91,Markarian:94,Gutkin:99}. It is worth mentioning that, while generic convex billiards have positive topological entropy \cite{BessaDelMagnoLopesDiasGaivaoTorres24, Cheng2004400,ZHANG2017793}, this fact alone does not yield examples with positive metric entropy with respect to $m$.

%Our first main result characterizes precisely when Lyapunov exponents vanish for random billiards, where $f$ denotes the billiard map on $\Tt^2$, $f_x$ its additive perturbations, and $F_\omega^n$ the associated random compositions. Additive perturbations $f_x$ of the billiard map $f$ admit a clear geometric interpretation. Given $x=(\bar s,\bar r)$ and an initial collision $(s_0,r_0)$, one first computes the next collision $(s_1,r_1)=f(s_0,r_0)$, and then translates the collision point by $x$ in the phase space, shifting the arc-length coordinate by $\bar s$ modulo $L$ and the $r$-coordinate by $\bar r$ modulo $2$.

%\begin{theorem}
%\label{thm:main}
%Let $D$ be a convex billiard table on a surface $S$ of constant curvature. Assume that the common distribution $\nu$ of the random variables $\{\omega_n\}_{n \ge 0}$ is absolutely continuous with respect to $m$. Then the Lyapunov exponents of the random compositions $F_\omega^n$ are non-zero almost everywhere if $D$ is not a geodesic disk, and vanish almost everywhere otherwise.
%\end{theorem}

%This theorem reflects the exceptional symmetry of circular billiards on surfaces of constant curvature. For such billiards, the phase space of the map $f$ is foliated by non-contractible invariant circles \cite{Bialy:1993ty,dosSantos:2017}, a property often referred to as \emph{total integrability}. When $f$ is replaced by an additive perturbation $f_x$, this foliation remains globally invariant: although individual circles may no longer be invariant under $f_x$, each circle of the foliation is mapped by $f_x$ onto another circle. As a consequence, the associated random dynamical system exhibits no exponential growth of derivatives, and hence has vanishing Lyapunov exponents almost everywhere.

%Theorem~\ref{thm:main} significantly extends the 
%result obtained in~\cite{Nguyen2021}, where it was shown that random additive perturbations of billiards in non-circular elliptical tables in the Euclidean plane have non-zero Lyapunov exponents almost everywhere.


%The second class of conservative dynamical systems considered in this paper consists of generalized standard maps. A \emph{generalized standard map} is a map $f_V\colon\Tt^2\to\Tt^2$ defined by
%\[
%f_V(y_1,y_2) = \left(y_1+y_2+V(y_1),\, y_2+V(y_1)\right) \bmod 1,
%\]
%where $V \colon\Tt^1\to\Rr$ is a $C^1$ function. For $V(y_1)=K\sin(2\pi y_1)$, one obtains the standard map introduced by Chirikov~\cite{CHIRIKOV1979263}; the case $K=0$ is completely integrable.

%Generalized standard maps arise naturally in Hamiltonian dynamics as models of periodically kicked systems and play a central role in the study of the transition from integrability to chaos. However, proving the existence of non-zero Lyapunov exponents for non-integrable standard maps and their generalizations remains notoriously difficult.

%Our second main result establishes a rigidity theorem for random additive perturbations of standard maps, analogous to Theorem~\ref{thm:main}. In this setting, the conclusion holds under weaker assumptions on the common distribution of the random variables $(\omega_n)_{n\ge0}$.

%We call a Borel probability measure $\nu$ on $\Tt^2$ a \emph{vertical line measure} if $\nu=\delta_0 \times \chi$, where $\delta_0$ is the Dirac measure at $0\in\Tt^1$, $\chi=h\,m_1$ and $m_1$ denotes the normalized Lebesgue measure on the factor $\Tt^1$. 

%\begin{theorem}\label{thm:standard}
%Let $f_V\colon\Tt^2\to\Tt^2$ be a generalized standard map. Assume that the common distribution $\nu$ of the random variables $\{\omega_n\}_{n \ge 0}$ is either absolutely continuous with respect to $m$ or a vertical line measure $\delta_0 \times \chi$. Then the Lyapunov exponents of the random compositions $F_\omega^n$ are non-zero almost everywhere if $V$ is not constant, and vanish almost everywhere otherwise.
%\end{theorem}

%For random standard maps, Blumenthal, Xue, and Young obtained quantitative lower bounds on the maximal Lyapunov exponent under sufficiently strong random additive perturbations~\cite{BlumenthalXueYoung}. In contrast, the present work does not rely on large randomness and instead provides a sharp characterization of the vanishing of Lyapunov exponents.

%Taken together, Theorems~\ref{thm:main} and~\ref{thm:standard} show that, for random conservative perturbations of billiards and standard maps, vanishing Lyapunov exponents occur if and only if the underlying deterministic system is totally integrable.

%The conceptual core of the paper is a rigidity result for random dynamical systems, which we call an invariance principle, following the terminology of Avila--Viana~\cite{Avila:2010aa}. Roughly speaking, we show that if the Lyapunov exponents of a random system coincide almost everywhere, then there exists a non-random invariant structure on the projective tangent bundle. This is stated precisely in Theorem~\ref{thm:G} and is derived by applying a theorem of Ledrappier~\cite{Ledrappier:1986ue} to the setting of random maps considered in this paper. The argument is rooted in Furstenberg's theory of random matrix products~\cite{Furstenberg:1963uq} and its subsequent extensions to stationary sequences of matrices by Ledrappier~\cite{Ledrappier:1986ue} and to random diffeomorphisms on manifolds by Carverhill~\cite{carverhill1987furstenberg} and Baxendale~\cite{Baxendale:1989aa}.

%The paper is organized as follows. Section~\ref{sec:random maps} establishes a general invariance principle for random piecewise smooth maps and derives several consequences needed later in the paper. Section~\ref{se:ergodicity} proves the unique ergodicity of random additive perturbations and the equidistribution of their random orbits in two general settings. These settings cover the random additive perturbations of billiard maps and generalized standard maps treated in the final sections. Section~\ref{sec:billiards} provides background on convex billiards on surfaces of constant curvature. The proof of Theorem~\ref{thm:main} is given in Section~\ref{ss:random_billiards}. Finally, Section~\ref{sec:random_standard} is devoted to random generalized standard maps and contains the proof of Theorem~\ref{thm:standard}.

%---
%
%\added{Old version of the Introduction}
%
%Let $S$ be a Riemannian surface $S$ of constant curvature, and let $D \subset S$ be a strictly convex domain whose boundary $\partial D$ is a $C^2$ simple closed curve of positive geodesic curvature at every point. We refer to such a domain $D$ as a \emph{convex table}.  
%
%A \emph{billiard} in $D$ is the dynamical system consisting of a point particle that moves along geodesics of $S$ within $D$ and reflects elastically at the boundary $\partial D$ so that the angle of reflection equals the angle of incidence. The motion of the particle can be described by the \emph{billiard map} $\Phi \colon V \to V$, where $V=S^1 \times [-1,1]$ represents the space of collisions of the particle with $\partial D$. More precisely, let $L$ denote the length of $\partial D$. Each pair $(s,r) \in V$ uniquely identifies a collision of the particle with $\partial D$: the coordinate $s \in \Rr / (L\Zz) \cong S^1$ is the arc-length of the point of collision of the particle off $\partial D$, while $r=-\cos \theta \in [-1,1]$, where $\theta \in [0,\pi]$ is the reflection angle of the particle measured with respect to the tangent of $\partial D$ at $s$. The billiard map $\Phi$ assigns to each collision $(s,r) \in V$ the next collision along the particle's trajectory. 
%%The precise definition of $\Phi$ is given in Section~\ref{}. 
%It is well known that $\Phi$ is a twist map preserving the Lebesgue measure $dsdr$ on $V$. This implies, in particular, that the sum of the Lyapunov exponents of $\Phi$ is equal to zero. 
%
%Numerical experiments suggest that, for a generic convex table $D$, the set $V$ splits into two invariant subsets of positive Lebesgue measure:  on one, all Lyapunov exponents of $\Phi$ vanish, while on the other, one Lyapunov exponent is positive and the other is negative. If this property were rigorously established, then, by Pesin's formula, $\Phi$ would have positive metric entropy with respect to the Lebesgue measure. However, no billiard with strictly convex table and positive metric entropy is currently known. Notably, all known examples of billiards with positive metric entropy have tables that are neither strictly convex nor have $C^2$ boundaries: for billiards in the Euclidean plane, see \cite{Bunimovich:79, Wojtkowski:86, Donnay:91, Markarian:94}; for billiards on the sphere and the hyperbolic plane, see \cite{Gutkin:99}.  We also remark that although the topological entropy of $\Phi$ is positive for generic convex tables \cite{}, this property alone does not imply positive metric entropy. 
%
%Given the difficulty of establishing the positivity of metric entropy for billiards with convex tables, we adopt an alternative approach: rather than studying the deterministic map $\Phi$, we consider a suitable random perturbation of $\Phi$. More specifically, we consider an additive random perturbation of $\Phi$,  constructed by randomly selecting, for every initial condition $(s,r)$, a point in a small neighborhood of $\Phi(s,r)$ in $V$. 
%
%To define this perturbation precisely, we first identify the collision space $V = S^1 \times [-1,1]$ with $\Tt^2$ by glueing together the point $-1$ and $1$ of the interval. This corresponds to identify the reflection angles $\theta=0$ and $\theta=\pi$. Let  $T \colon \Tt^2 \to  \Tt^2$ be the map induced by $\Phi$ on $\Tt^2$; full details of this construction are provided in Section~\ref{sec:billiards}.  We then introduce a family $\{T_x\}_{x \in \Tt^2}$ of deterministic perturbations of $T$ defined by 
%\[
%T_x(y) = T(y) + x \quad \text{for all } y \in \Tt^2. 
%\]
%Geometrically, this means the following procedure: given a perturbation parameter $x=(\bar{s},\bar{r})$ and an initial collision $(s_0,r_0)$, we compute the next collision $(s_1,r_1)=T(s_0,r_0)$, then shift the collision point from $s_1$ to $s_1+\bar{s} \bmod L$ and the $r$-coordinate from $r_1$ to $r_1+\bar{r} \bmod 2$. 
%
%To introduce randomness, we consider an i.i.d. sequence of real random variables  $\{X_n\}_{n \ge 0}$ taking values in a small neighborhood of $0 \in \Tt^2$. Each $T_{X_n}$ is then a \emph{random additive perturbation} of $T$. The random dynamical system generated by the maps $T_{X_n}$ is the sequence of compositions
%\[
%\mathsf{T}^{n} = T_{X_{n-1}} \circ \cdots \circ T_{X_0} \quad \text{for all } n \ge 1,
%\]
%which may be viewed as the random analogue of the deterministic iterates $T^n$. We call the sequence $\{\mathsf{T}^n\}$ a 
%%\added[comment=a more informal terminology: \emph{random perturbation of the billiard} with table $D$]
%\emph{random additive perturbation} of the billiard map $\Phi$. Given a point $x_0 \in \Tt^2$, its random orbit is the random process $\{Y_n\}$ defined recursively by
%\[
%Y_{n+1} = T_{X_n}(Y_n), \quad Y_0=x_0.
%\]
%This process represents the counterpart of a deterministic billiard orbit of $x_0$.
%
%In this paper, we address the following question: for which convex table $D$ of a surface $S$ of constant curvature, does the corresponding randomly dynamical system $\{\mathsf{T}^n\}$ exhibit non-zero Lyapunov exponents almost everywhere? 
%
%This problem was first investigated by Nguyen in the setting of billiards with convex tables in the Euclidean plane~\cite{Nguyen2021}, where it was shown that the Lyapunov exponents are non-zero almost everywhere for elliptical tables that are not circular. 
%
%In the present work, we significantly extend Nguyen's results by proving that, under suitable conditions on the common distribution of the random variables $X_n$ (see Section~\ref{ss:random_billiards}), that the Lyapunov exponents of $\{\mathsf{T}^n\}$ vanish if and only if the convex table $D$ is a geodesic disk in a surface of constant curvature. In all other cases, a random additive perturbation yields non-zero Lyapunov exponents almost everywhere. 
%
%\begin{theorem}
%\label{thm:main}
%Let $D$ be a convex billiard table on a surface $S$ of constant curvature, and assume that the common distribution of the random variables $X_n$ is absolutely continuous with respect to the Lebesgue measure on $\mathbb{T}^2$, with support containing a nonempty open set. Then the Lyapunov exponents of the random dynamical system $\{\mathsf{T}^n\}$ vanish if and only if $D$ is a geodesic disk.
%\end{theorem}
%
%%\begin{theorem}
%%\label{thm: main}
%%Let $D$ be a convex table of a surface $S$ of constant curvature. Suppose that the common distribution of the random variables $X_n$ is absolutely continuous with respect to the volume measure of $\Tt^2$ and contains an open set. 
%%(see Section~\ref{ss:random_billiards}), 
%%Then the random dynamical system $\{\mathsf{T}^n\}$ has  zero Lyapunov exponents if and only if $D$ is a geodesic disk.
%%\end{theorem}
%
%This result reflects the exceptional symmetry of circular billiards on surfaces of constant curvature. For these billiards, the collision space $V$ is foliated by non-contractible circles invariant under $T$~\cite{Bialy,dos Santos}. When $T$ is replaced by any additive perturbation $T_x$, this foliation remains globally invariant: although individual circles are no longer invariant under $T_x$, the map $T_x$ sends each circle to another circle of the foliation. 
%%Our result also highlights a possible connection with the Birkhoff conjecture, which posits that ellipses are the only integrable strictly convex billiard tables.
%
%%Our proof consists of two steps. 
%The main step in the proof of Theorem~\ref{thm:main} consists in establishing a rigidity result that applies to a broad class of random dynamical systems on a smooth manifold $M$, including the random perturbations of billiard maps considered here. Informally, we prove that if all Lyapunov exponents of such a system coincide, then there exists a non-random invariant structure in the projective tangent bundle of $M$ (see Theorem~\ref{thm: G}). To achieve this, we adapt a theorem of Ledrappier to our setting (see Theorem~\ref{thm: Ledrappier}).
%
%This rigidity result belongs to a class of theorems whose prototype is Furstenberg's classical theorem on products of i.i.d. random matrices in $\operatorname{GL}(d,\mathbb{R})$~\cite{Furstenberg:1963uq}. Furstenberg's theorem asserts that, if all Lyapunov exponents coincide, then there exists a probability measure on the projective space $\mathbb{P}^{d-1}$ that is invariant under the projective action of every matrix in the support of the common distribution. This conclusion provides a practical criterion for the simplicity of the Lyapunov spectrum: to show that the Lyapunov exponents are distinct, it suffices to rule out the existence of such an invariant probability measure on $\mathbb{P}^{d-1}$.
%
%Extensions of Furstenberg's theorem have been obtained in a various contexts. For linear cocycles over stationary Markov chains, key results were obtained by Royer~\cite{Royer:1980vh}, Guivarc'h~\cite{Guivarch:79} and Ledrappier~\cite{Ledrappier:1986ue}. For random diffeomorphisms on manifolds, analogous results are due to Carverhill~\cite{carverhill1987furstenberg} and %Crauel~\cite{Crauel:90}, 
%Baxendale~\cite{Baxendale:1989aa}. In deterministic settings, Bonatti, G\'omez-Mont and Viana~\cite{Bonatti:03} obtained Furstenberg-type criteria for H\"older-continuous linear cocycles over hyperbolic systems. The most general result in this direction is the Invariance Principle of Avila and Viana~\cite{Avila:2010aa}.
%
%To complete the proof of Theorem~\ref{thm:main}, we proceed as follows. Applying Theorem~\ref{thm: G} to the random dynamical system ${\mathsf{T}^n}$, we deduce that if all Lyapunov exponents coincide, there exists a probability measure on the projective tangent bundle of $\Tt^2$ that is invariant under the projective action of the derivative $DT$. The existence of such an invariant measure yields a foliation of the collision space $V$ by non-contractible circles invariant under the billiard map $\Phi$, as occurs in the case of circular billiards. The final step consists in showing that this foliation forces the table $D$ to be a geodesic disk. 
%
%%\begin{theorem}
%%Let $g \colon S^1 \times [a,b] \to S^1 \times [a,b]$ be a twist diffeomorphism that preserves the area and fixes at least one of the two boundary circles $S^1 \times \{a\}$ or $S^1 \times \{b\}$. Then any random additive pertubation of $g$ has zero Lyapunov exponents if and only if $g$ is the integrable twist map  $g(x,y) = (x+y \bmod 1,y)$ .
%%\end{theorem}
%
%The paper is organized as follows. Section~\ref{sec:random maps}  reviews the representation of random dynamical systems as skew-products over Bernoulli shifts. We introduce the standing assumptions on the class of random maps under consideration and define the associated extremal Lyapunov exponents. We then recall Ledrappier's theorem, which generalizes Furstenberg's classical result on the vanishing of Lyapunov exponents for products of independent unimodular matrices~\cite{Furstenberg:1963uq} to linear cocycles over Markov chains. Using Ledrappier's theorem, we extend Furstenberg's result to the class of random maps introduced earlier. We further specialize this result to random additive perturbations of toral maps and derive several corollaries. The section concludes with a summary of results on the ergodic properties of the Markov chains and the skew-products associated with these random dynamical systems. 
%
%In Section~\ref{sec:billiards}, we provide background on convex billiards on surfaces of constant curvature. We describe the billiard map in different coordinate systems and recall that billiards in geodesic disks are integrable systems. 
%
%Section~\ref{ss:random_billiards} introduces billiard maps with random additive perturbations. We verify that these systems satisfy the assumptions of the results established in Section~\ref{sec:random maps}, and we complete the proof of Theorem~\ref{thm: main} by showing that circular billiards are the only ones for which random additive perturbations lead to vanishing Lyapunov exponents. 
%%This is proved for both absolutely continuous perturbations and singular perturbations along the $r$-coordinate.
%
%%Finally, 
%In Section~\ref{sec:final} we discuss extensions of Theorem~\ref{thm:main} to magnetic billiards.
%% and additional applications of our general results from Section~\ref{sec:random maps} to \added[comment=and non-twist map since the twist property is used only in Proposition~\ref{prop:singular}]{twist maps and in particular to the standard map}.
%
%Finally, in Section~\ref{sec:random maps}, we obtained results analogous to Theorem~\ref{thm:main} for random standard maps and some of their extension under various assumptions on the distribution of the i.i.d. random variables generating randomness.
%




%%%%%%%%%%%%%%%%%%%%%%%%%%%%%%%%%%%%%%%%%%%%%%%%%%%%%%%%%%%%%
%\section{Outer billiards}
%\label{sec:Outer billiards}
%\input{outer_billiards.tex}

%%%%%%%%%%%%%%%%%%%%%%%%%%%%%%%%%%%%%%%%%%%%%%%%%%%%%%%%%%%%%
\section{An invariant principle for random maps}
%\section{Lyapunov exponents of random dynamical systems}
\label{sec:random maps}
\subsection{Random dynamical systems} 
We model a random dynamical system as a skew-product over a shift transformation. The definition of a random dynamical system provided below is not the most general. For more general definitions, see \cite{Arnold:98,Kifer:86}.
%it is specifically tailored to describe random additive perturbations of a deterministic map 

%--- Remark: the compactness of M is not necessary; we just need M to be complete and separable. The easiest way to obtain this is to require M be compact. Connectedness is not needed.

Throughout the paper, $\Nn$ denotes the natural numbers, including zero. Let $M$ be a complete smooth Riemannian manifold of dimension $d$, equipped
with the Borel $\sigma$-algebra $\BB$. We denote by $\operatorname{vol}$ the Riemannian volume measure on $M$.

Let $(X, \XX, \nu)$ be a probability space, where $X$ is a separable complete metric space, $\XX$ is its Borel $\sigma$-algebra, and $\nu$ is a probability measure on $\XX$. Denote by $(\Omega,\Sigma,\rho_\nu)$ the probability space given by the product of countably many copies of $(X, \XX, \nu)$. Namely, $\Omega$ is the space of sequences 
\[
\Omega = X^{\Nn} = \{(\omega_n)_{n \in \Nn} \colon \omega_n \in X \text{ for all } n \in  \Nn\},
\] 
equipped with the product $\sigma$-algebra $\Sigma = \XX^{\Nn}$ and the product measure $\rho_\nu = \nu^{\Nn}$. The shift map $\sigma \colon \Omega \to \Omega$ defined by
\[(\sigma(\omega))_n = \omega_{n+1} \quad \text{for every } \omega \in \Omega,\] 
is measurable and preserves the probability measure $\rho_\nu$.

Since $X$ and $M$ are separable complete metric spaces, the same is true of $X\times M$, $\Omega$ and $\Omega\times M$. All these spaces are endowed
with their Borel $\sigma$-algebras.
%, which coincide with the product $\sigma$-algebra of the corresponding factor spaces \cite[Lemma 1.2]{Kallenberg:97}.


%Consider a measurable map $f \colon X \times M \to M$.

%Define $\tilde{f} \colon \Omega \times M \to M$ by 
%\[
%\tilde{f}(\omega,y) = f(\omega_0,y) \quad \text{for every } (\omega,y) \in \Omega \times M.
%\]
%The map $\tilde{f}$ retains the properties of $f$: it is measurable, and for every $\omega \in \Omega$, the map $\tilde{f}(\omega,\cdot) \colon M \to M$ is a homeomorphism and its restriction to $N$ is a $C^1$ embedding. 


%where $\pi \colon \Omega \times M \to X \times M$ is the projection given by $\pi(\omega,y) = (\omega_0,y)$ for every $(\omega,y) \in \Omega \times M$. 

Let $f \colon X \times M \to M$ be a measurable map. For every $x \in X$, define 
\[
f_x(\cdot)=f(x,\cdot) \colon M \to M.
\]

\begin{defn}
We call the skew-product $F \colon \Omega \times M \to \Omega \times M$ defined by
\[
F(\omega,y)=(\sigma(\omega),f_{\omega_0}(y)) \quad  \text{for every } (\omega,y) \in \Omega \times M,
\]  
the \emph{random dynamical system on $M$ generated by $f$ and $\nu$}.
\end{defn}

Note that $F$ depends on $\omega$ only through its 0th component $\omega_0$. For every $n \in \Nn$ and every $ \omega \in \Omega$, define $F^{n}_\omega \colon M \to M$ as follows: 
\[
F^{n}_\omega = 
\begin{cases}
\operatorname{id}_M & \text{if } n=0, \\
f_{\omega_{n-1}} \circ \cdots \circ f_{\omega_0}  & \text{if } n \ge 1.
%f^{-1}_{\omega_{n}} \circ \cdots \circ f^{-1}_{\omega_{-1}} (y) & \text{if } n \le -1. 
\end{cases}
\]
The iterates of $F$ can be written as 
\[
F^{n}(\omega,y)=(\sigma^n(\omega),F^{n}_{\omega}(y)) \quad  \text{for every } (\omega,y) \in \Omega \times M.
\]

\begin{defn}
Let $y \in M$. For every $\omega \in \Omega$, the sequence $\{F^n_{\omega}(y)\}_{n \ge 0}$ is called the \emph{random orbit of $y$ associated with $\omega$}.
\end{defn}

We now formulate the standing assumptions on $f$ and $\nu$ that will be assumed throughout this paper.

\begin{assumption}\label{as:one}
There exists a Borel probability measure $m$ on $M$ that is $f_x$-invariant for every $x\in X$. Moreover, for every $x\in X$, there exists an open set $N_x\subset M$ with
$m(N_x)=1$ such that
\[
    f_x|_{N_x}\colon N_x\to f_x(N_x)\subset M
\]
is a $C^1$ diffeomorphism. Finally, the set
\[
    \NN:=\{(\omega,y)\in\Omega\times M:\ y\in N_{\omega_0}\}
\]
is measurable in $\Omega\times M$.
\end{assumption}

%\begin{assumption}\label{as:one}
%There exists a Borel probability measure $m$ on $M$ absolutely continuous with respect to $\operatorname{vol}$ such that for every $x\in X$,
%\begin{enumerate}
%    \item $m$ is invariant under $f_x$,
%    \item there exists an open set $N_x\subset M$ with
%    $m(N_x)=1$ such that
%    $f_x|_{N_x}\colon N_x \to f_x(N_x) \subset M$ is a $C^1$ embedding.
%\end{enumerate}
%Moreover, $\NN:=\{(\omega,y):\omega \in \Omega, \,y\in N_{\omega_0}\}$ is a measurable subset of $\Omega \times M$.
%\end{assumption}

\begin{remark}
Assumption~\ref{as:one} allows the maps $f_x$ to have singular points, as in the case of billiard maps. We call a point $y\in M$ singular for $f_x$ if $f_x$ is not a $C^1$ local diffeomorphism at $y$. Since
$f_x|_{N_x}$ is a $C^1$ diffeomorphism, every singular point of $f_x$ belongs to $M\setminus N_x$. Moreover, the invariance of $m$ under $f_x$ for every $x\in X$ implies that $\rho_\nu\times m$ is invariant under the skew product $F$. Finally, since $m(N_x)=1$ for every $x\in X$, Fubini's theorem gives $(\rho_\nu\times m)(\NN)=1$.
\end{remark}

We fix a global measurable trivialization $\Psi\colon TM\to M\times\Rr^d$ of the tangent bundle such that for each $y\in M$, the restriction
\[
\Psi_y:=\Psi|_{T_yM}\colon T_yM\to\Rr^d
\]
is a linear isometry from $T_yM$, endowed with the inner product induced by the Riemannian metric, onto $\Rr^d$, endowed with its standard Euclidean inner product; see \cite[Lemma~4.2.4]{Arnold:98}.

Define $A\colon \Omega\times M\to \operatorname{GL}(d,\Rr)$ by
\[
A(\omega,y)
=
\begin{cases}
\Psi_{f_{\omega_0}(y)}\circ Df_{\omega_0}(y)\circ \Psi_y^{-1}
& \text{if } (\omega,y)\in \NN,\\[2mm]
I
& \text{if } (\omega,y)\in (\Omega\times M)\setminus \NN,
\end{cases}
\]
where $I$ denotes the identity matrix in $\operatorname{GL}(d,\Rr)$.

On $\NN$, the matrix $A(\omega,y)$ represents the map $Df_{\omega_0}(y)$ with respect to the isometric trivialization $\Psi$. Therefore
\[
\|A(\omega,y)\|
=
\|Df_{\omega_0}(y)\|_{y,f_{\omega_0}(y)}
\qquad\text{for all }(\omega,y)\in\NN,
\]
where the norm on the left-hand side is the Euclidean operator norm on $\Rr^d$, while the norm on the right-hand side is the Riemannian operator norm from $T_yM$ to $T_{f_{\omega_0}(y)}M$.

Since $\NN$ is measurable, $\Psi$ is measurable, and
$(\omega,y)\mapsto Df_{\omega_0}(y)$ is measurable on $\NN$, the map $A$ is measurable on $\Omega\times M$. We set $A(\omega,y)=I$ on
$(\Omega\times M)\setminus\NN$ only to obtain a measurable map defined everywhere.

Set $A^0(\omega,y)=I$. For $n\ge 1$, define
\[
A^n(\omega,y)
=
A(F^{n-1}(\omega,y))\cdots A(F(\omega,y))A(\omega,y).
\]

Note that the choice of $A$ on $(\Omega\times M)\setminus\NN$ does not affect
$A^n(\omega,y)$ on the full $\rho_\nu\times m$-measure set of points whose forward $F$-orbit stays in $\NN$.

\begin{defn}
The maximal and minimal Lyapunov exponents of $F$ at $(\omega,y) \in \Omega \times M$ are defined by
\[
\lambda_F^+(\omega,y) = \limsup_{n\to+\infty} \frac1n\log\|A^n(\omega,y)\|
\]
and
\[
\lambda_F^-(\omega,y)
= -\limsup_{n\to+\infty}
\frac1n\log\|(A^n(\omega,y))^{-1}\|.
\]
\end{defn}

\begin{assumption}\label{as:three}
$\log^+\|A\|,\ \log^+\|A^{-1}\|
    \in L^1(\rho_\nu\times m)$.
\end{assumption}

\begin{remark}
By Assumption~\ref{as:three} and the Subadditive Ergodic Theorem, the limsup in the definitions of $\lambda_F^+(\omega,y)$ and $\lambda_F^-(\omega,y)$ can be replaced by the limit for $\rho_\nu\times m$-a.e. $(\omega,y)\in\Omega\times M$. Since the Lyapunov exponent are $F$-invariant, ergodicity of
$\rho_\nu\times m$ with respect to $F$ implies that both exponents are constant $\rho_\nu\times m$-a.e.
\end{remark}

Since $(\rho_{\nu} \times m)(\NN)=1$, the set of points $(\omega,y)$ such that
\[
F^n(\omega,y) \in \NN
\qquad\text{for all } n\ge 0
\]
has full $\rho_\nu\times m$-measure. On this set, $A^n(\omega,y)$ is the matrix representation of $DF_\omega^n(y)$ with respect to $\Psi$. Hence for $\rho_\nu\times m$-a.e. $(\omega,y)$, 
\[
\lambda_F^+(\omega,y)
=
\lim_{n\to+\infty}
\frac1n\log\|DF_\omega^n(y)\|_{y,F_\omega^n(y)}
\]
and
\[
\lambda_F^-(\omega,y)
=
-\lim_{n\to+\infty}
\frac1n\log\|(DF_\omega^n(y))^{-1}\|_{F_\omega^n(y),y}.
\]

\subsection{Random additive perturbations of toral maps}
\label{random additive}

We now specialize to random dynamical systems for which both the base space
and the fiber are the $d$-dimensional torus
\[
X=M=\Tt^d=\Rr^d/\Zz^d.
\]
We endow $\Tt^d$ with its standard flat metric and denote by $m$ its normalized Haar measure.

For brevity, throughout this paper the term \emph{toral map} will be used in the specialized sense specified in the following definition. In particular, it will always refer to a map preserving $m$ and satisfying the stated regularity and integrability assumptions.

\begin{defn}
\label{def:toral_map}
A measurable map $g\colon\Tt^d\to\Tt^d$ is called a \emph{toral map} if it satisfies the following properties:
\begin{enumerate}
\item $g$ preserves the Haar measure $m$, 
\item there exists an open set $N\subset\Tt^d$ with $m(N)=1$ such that $g|_N\colon N\to g(N) \subset \Tt^d$
is a $C^1$ diffeomorphism,
\item the functions $1_N(y)\log^+\|(Dg(y))^{\pm 1}\|$  belong to $L^1(m)$, where they are understood to vanish on $\Tt^d\setminus N$ and the norm is the operator norm induced by the
standard flat metric on $\Tt^d$. 
\end{enumerate}
\end{defn}

\begin{defn}
Let $g\colon\Tt^d\to\Tt^d$ be a toral. A map
$f\colon\Tt^d\times\Tt^d\to\Tt^d$ of the form
\[
f(x,y)=\tau_x\circ g(y)=g(y)+x,
\qquad x,y\in\Tt^d,
\]
where $\tau_x$ denotes translation by $x$ on $\Tt^d$, is called an
\emph{additive perturbation} of $g$. The random dynamical system generated by such a map $f$ and by a probability measure $\nu$ on $\Tt^d$ is called a \emph{$\nu$-random additive
perturbation of $g$}. In this setting, $\Omega=(\Tt^d)^{\Nn}$ and $\rho_\nu=\nu^{\Nn}$.
\end{defn}

%\begin{defn}
%A map $f\colon \Tt^d\times\Tt^d\to\Tt^d$ satisfying the preceding assumptions is called an \emph{additive perturbation} of the toral map $g$. The random dynamical system generated by such an $f$ and by a probability measure $\nu$ on $\Tt^d$ is called a \emph{$\nu$-random additive perturbation of $g$}.
%\end{defn}

%In this setting, $\nu$ is a probability measure on $\Tt^d$, and $\rho_\nu$ is the corresponding product measure on $\Omega=(\Tt^d)^{\Nn}$.

%\begin{defn}
%A map $f\colon \Tt^d\times\Tt^d\to\Tt^d$ satisfying the preceding
%assumptions is called an \emph{additive perturbation} of the toral map $g$.
%The random dynamical system generated by such an $f$ and by a probability
%measure $\nu$ on $\Tt^d$ is called a \emph{$\nu$-random additive
%perturbation of $g$}.
%\end{defn}

Random additive perturbations toral maps provide examples of systems satisfying Assumptions~\ref{as:one} and~\ref{as:three}.

\begin{lemma}
\label{le:add-pert}
Every $\nu$-random additive perturbation $F$ of a toral map satisfies Assumptions~\ref{as:one} and~\ref{as:three}. Moreover, for such a system, 
the condition $\lambda_F^-(\omega,y)=\lambda_F^+(\omega,y)$ for $\rho_\nu\times m$-a.e. $(\omega,y)$ is equivalent to the condition $\lambda_F^+(\omega,y)=0$ for $\rho_\nu\times m$-a.e. $(\omega,y)$.
\end{lemma}

\begin{proof}
Translations of $\Tt^d$ are diffeomorphisms preserving the Haar measure $m$. Hence, for every $x\in\Tt^d$, the map $f_x=\tau_x\circ g$ preserves $m$, and $f_x|_N\colon N\to f_x(N)=\tau_x(g(N))$ is a $C^1$ diffeomorphism. Moreover, taking $N_x=N$ for every $x\in\Tt^d$, we have
$\NN=\Omega\times N$, which is measurable. Thus $F$ satisfies Assumption~\ref{as:one}.

Since the differential of a translation is equal to the identity in the standard flat trivialization, we have 
\[
Df_x(y)=Dg(y) \qquad \text{for all } (x,y) \in\Tt^d \times N.
\]
Thus $\|(A(\omega,y))^{\pm 1}\|=\|(Dg(y))^{\pm 1}\|$ on $\NN$. Since $A(\omega,y)=I$ on $(\Omega \times \Tt^d) \setminus \NN$, the integrability assumptions on $(Dg)^{\pm 1}$ imply that $\log^+ \|A^{\pm 1}\| \in L^{1}(\rho_{\nu} \times m)$. Therefore $F$ satisfies Assumption~\ref{as:three}.


Furthermore, $|\det A^n(\omega,y)|=1$
for every $n\geq1$ and for $\rho_\nu\times m$-a.e.\ $(\omega,y)$. By Oseledets' Theorem,
\[
\lim_{n\to\infty}\frac1n\log|\det A^n(\omega,y)|
\]
is equal to the sum of the Lyapunov exponents of $F$, counted with multiplicity. Since the left-hand side is zero, this sum vanishes for $\rho_\nu\times m$-a.e.\ $(\omega,y)$.

If $\lambda_F^-=\lambda_F^+$ almost everywhere, then all Lyapunov exponents coincide almost everywhere, since every exponent lies between the minimal
and maximal ones. Their sum being zero, they must all vanish, and in particular $\lambda_F^+=0$ almost everywhere. The converse is immediate
from $\lambda_F^- \le \lambda_F^+$ and the vanishing of the sum of the Lyapunov exponents.
\end{proof}


%\begin{remark}
%\label{remark: one}

%Under the assumptions on $g$, any additive perturbation $F$ of $g$ satisfies Assumptions~\ref{as:one} and~\ref{as:three}. By Oseledets' Theorem~\cite{Viana:2014uo}, the limit
%\[
%\lim_{n\to+\infty}\frac{1}{n}\log\left|\det DF^n_\omega(y)\right|
%\]
%exists for $(\rho_\nu\times m)$-almost every $(\omega,y)\in\Omega\times\mathbb T^d$ and is equal to the sum of the Lyapunov exponents of $F$. Since $g$ is volume-preserving, we have
%\[
%Df_x(y)=Dg(y)\in \SL'(d,\mathbb R)
%\quad\text{for all }(x,y)\in\mathbb T^d\times\mathbb T^d.
%\]
%It follows that the limit above is equal to $0$. Consequently, for any additive random perturbation $F$, the condition $\lambda_F^+=\lambda_F^-$ a.e. on $\Omega\times\mathbb T^d$ in Theorem~\ref{thm: G} is equivalent to $
%\lambda_F^+=0$ a.e. on $\Omega\times\mathbb T^d$.
%\end{remark}

\subsection{Invariance principle}
%\subsection{Rigidity results for Lyapunov exponents}
%\subsection{Equality of the Lyapunov exponents}
In this section, we first recall a theorem of Ledrappier that generalizes Furstenberg's classical result on zero Lyapunov exponents for products of independent unimodular matrices~\cite{Furstenberg:1963uq} to linear cocycles over Markov chains. Similar results have been obtained in \cite{Royer:1980vh,Virtser:80,Guivarch:79}. Then we apply Ledrappier's theorem to random maps.

\subsubsection{Ledrappier's Invariance Theorem}

Let $Z$ be a separable complete metric space with its Borel $\sigma$-algebra $\ZZ$. Consider a Markov chain $\{Z_n\}_{n \ge \Nn}$ with state space $(Z,\ZZ)$, transition probability $\{P_z \colon z \in Z\}$ and stationary probability measure $\zeta$. 

Let $(\Omega',\Sigma')$ be the two-sided product space $\Omega'=Z^{\Zz}$ endowed with the product $\sigma$-algebra $\Sigma' := \ZZ^{\Zz}$. For each $n \in \Zz$, let $\pi_n \colon \Omega' \to Z$ be the projection given by $\pi_n(\omega')=\omega'_n$ for all $\omega'=(\omega'_n)_{n \in \Zz} \in \Omega'$. Since  $Z$ is a separable complete metric space, there exists a probability measure $\Pp_{\zeta}$ on $\Sigma'$ such that the coordinate process $\{\pi_n\}_{n \in \Zz}$ is a Markov chain on $(\Omega',\Sigma',\Pp_{\zeta})$ with transition probability $P_z$ and stationary probability measure $\zeta$. Moreover, the shift map $\sigma' \colon \Omega' \to \Omega'$ 
%on $\Omega'$ 
is measurable with respect to $\Sigma'$ and preserves $\Pp_{\zeta}$  (see \cite[Theorem 3.1.7 and Example 5.1.3]{Douc:2018wc}).
%\cite[Example 5.1.3]{Douc:2018wc} \added[comment=check the reference]{(see also \cite[Proposition 4.49]{Benaim_2022}).}

Let $A \colon Z \to \operatorname{GL}(d,\Rr)$ be a measurable map such $\log^+ \|A(\cdot)^{\pm 1}\|$ belong to $L^1(\zeta)$. Consider the linear cocycle over $(\Omega',\Sigma',\Pp_{\zeta},\sigma')$ generated by $A$, and denote by $\Lambda^+_A(\omega')$ and $\Lambda^-_A(\omega')$ its maximal and minimal Lyapunov exponents, respectively \cite{Ledrappier:1986ue}. 

By abuse of notation, we also denote by $A(z)$ the action induced by the matrix $A(z)$ on the projective space $\Pp^{d-1}$.

A family of probability measures $\{\eta_z : z \in Z\}$ on $\mathbb{P}^{d-1}$ is said to be \emph{measurable} if for every Borel set $B \subset \mathbb{P}^{d-1}$, the map
$z \mapsto \eta_z(B)$ is $\mathcal Z$-measurable. Since $\mathbb{P}^{d-1}$ is a compact metrizable space, this condition is equivalent to requiring that the map
$z \mapsto \eta_z$ is measurable as a map from $(Z,\mathcal Z)$ into the space of probability measures on $\mathbb{P}^{d-1}$ endowed with the Borel $\sigma$-algebra of the weak-* topology.

\begin{theorem}[{\cite[Corollary 2]{Ledrappier:1986ue}}]
\label{thm: Ledrappier}
Suppose that 
%$\Lambda^+_A(\omega') = \Lambda^-_A(\omega')$ for $\Pp_{\zeta}$-a.e. $\omega' \in \Omega'$. 
\[
\Lambda^+_A(\omega') = \Lambda^-_A(\omega') \quad \text{for } \Pp_{\zeta} \text{-a.e. } \omega' \in \Omega'.
\] 
Then there exists a measurable family $\{\eta_z : z \in Z\}$ of probability measures on $\mathbb{P}^{d-1}$ such that for $\zeta$-a.e. $z \in Z$, 
\[
A(z)_{*} \eta_z = \eta_{z_1} \quad \text{for } P_z  \text{-a.e. } z_1 \in Z.
\]
\end{theorem}

%The next remark explains why we can use Theorem \ref{thm: Ledrappier}.

%\begin{remark} 
%Observe that $Z:=X \times M$ is a complete separable metric space, since it is the product of $X$ and $M$ which are themselves complete separable metric spaces. Denote by $\ZZ$ the Borel $\sigma$-algebra of $Z$. Let $(\Omega',\Sigma')$ be the two-sided product space $\Omega'=Z^{\Zz}$ endowed with the product $\sigma$-algebra $\ZZ^{\Zz}$. For every $n \in \Zz$, let $\pi_n \colon \Omega' \to Z$ be the projection given by $\pi_n(\omega')=\omega'_n$ for all $\omega'=(\omega'_n)_{n \in \Zz} \in \Omega'$. There exists a probability measure $\Pp_{\mu}$ on $\Sigma'$ such that the coordinate process $\{\pi_n\}_{n \in \Zz}$ is a Markov chain on $(\Omega',\Sigma',\Pp_{\mu})$ with transition probability $P_z$ and initial  distribution $\mu$ (for example, see \cite[Theorem 3.1.7]{Douc:2018wc}). Moreover, the shift map $\sigma' \colon \Omega' \to \Omega'$ on $\Omega'$ is measurable with respect to $\Sigma'$ and preserves $(\Omega',\Sigma',\Pp_{\mu})$ \cite[Example 5.1.3]{Douc:2018wc} \added[comment=check the reference]{(see also \cite[Proposition 4.49]{Benaim_2022}).}
%\end{remark}

\subsubsection{Invariant principle for random maps}
%\subsubsection{Rigidity results for Lyapunov exponents of random maps}
%\subsubsection{A Furstenberg-type result for random maps}
%\subsubsection{Application of Ledrappier's Theorem to random maps}

We now establish a result (Theorem~\ref{thm:G}) that may be regarded as a generalization of Furstenberg’s theorem to random maps satisfying Assumptions~\ref{as:one} and~\ref{as:three}. Related results for random diffeomorphisms on manifolds were previously obtained by Carverhill~\cite{carverhill1987furstenberg} and Baxendale~\cite{Baxendale:1989aa}.

Denote by $PM = \bigsqcup_{y \in M} P_yM$ the projective bundle of the manifold $M$, where $P_yM$ is the projective space of $T_xM$. Let $\pi \colon PM \to M$ be the bundle projection. 

Any probability measure $\eta$ on $PM$ such that $\pi_* \eta = m$ admits a \emph{disintegration} with respect to $m$ (see \cite{Viana:2014uo}), i.e. a measurable family $\{\eta_y \colon y \in M\}$ of probability measures on $PM$ such that each $\eta_y$ is concentrated on the fiber $P_yM $ and   
\[
\eta(E) = \int_M \eta_y(E \cap P_y M) \, dm(y) \quad \text{for every measurable } E \subset PM.
\]  

We use the same notation $Df_x(y)$ for the differential of $f_x$ at $y$ and for the map it induces on the projective fiber $P_yM$.

%By abuse of notation, we also denote by $DT_x(y)$ the action induced by $DT_x(y)$ on the fibers of the projective bundle $PM$ of $M$.

\begin{theorem}
\label{thm:G}
Let $F \colon \Omega \times M \to \Omega \times M$ be a random dynamical system generated by $f$ and $\nu$ that satisfies Assumptions~\ref{as:one} and \ref{as:three}. Suppose that
\[
\lambda^+_F(\omega,y)=\lambda^-_F(\omega,y) \quad \text{for } \rho_\nu \times m \text{-a.e. } (\omega,y) \in \Omega \times M.
\] 
Then there exists a probability measure $\eta$ on $PM$ with $\pi_* \eta = m$ such that 
\[
(Df_x)_* \eta = \eta \quad \text{for } \nu \text{-a.e. } x \in X.
\]
In particular, if $\{\eta_y : y \in M\}$ denotes the disintegration of $\eta$ with respect to $m$, then 
\[
Df_x
(y)_* \eta_y = \eta_{f_x(y)} \quad \text{for } (\nu \times m) \text{-a.e. } (x,y) \in X \times M.
\]
\end{theorem}

\begin{proof}
The argument below follows closely the proof of \cite[Theorem 1]{carverhill1987furstenberg}.

Consider the Markov chain $\{(X_n,Y_n) \colon n \in \Nn\}$ on $X \times M$ and transition probabilities given by
\[
P_{(x,y)}(B \times C) = 	\nu(B) \delta_{f_x(y)}(C), \quad  B \in \XX, C \in \BB.
%\begin{cases}
%	\nu(B) \delta_{f_x(y)}(C) & \text{if } x \in X \text{ and } y \in N_x, \\
%	\nu(B) \delta_{y}(C) & \text{if } x \in X \text{ and } y \in M \setminus N_x, 
%\end{cases}
\]
for all $x \in X, y \in M$.
%\added{Old version:}
%\[
%P_{(x,y)}(B \times C) = \nu(B) \delta_{f_x(y)}(C)
%%= \nu(A) \chi_{B} \circ T_x(y)
%\]
%for all \deleted[comment=$y \in N_x$ otherwise $f_x(y)$ not defined]{$(x,y) \in X \times M$} and $B \in \XX, C \in \BB$. 
This means that $\{X_n\}$ is a sequence of i.i.d. random variables with values in $X$ and distribution $\nu$, whereas $\{Y_n\}$ is a Markov chain with state space $M$ satisfying the recursive relation   
\[
Y_{n+1}=f_{X_n}(Y_n).
\]
It is straightforward to check that $\mu:=\nu \times m$ is a stationary probability for $\{(X_n,Y_n)\}$.

Let $A\colon\Omega\times M\to\operatorname{GL}(d,\Rr)$ be the extension, defined after Assumption~\ref{as:one}, of the matrix representation of the derivative cocycle. Since $A(\omega,y)$ depends on $\omega$ only through its zeroth coordinate $\omega_0$, the map
\[
A' \colon X\times M\to\operatorname{GL}(d,\Rr),
\qquad
A'(x,y)=A(\omega,y),
\]
where $\omega\in\Omega$ is any sequence with $\omega_0=x$ is well defined.
%: for every $x \in X$ and every  $y \in N_x$, let $A(x,y)$ be the matrix representation of $Df_{x}(y)$ with respect to some global measurable trivialization of $TM$. 
%Note that $A$ is not defined for $y \in M \setminus N$. However, this is not an issue since $m(M \setminus N)=0$.

Under the hypotheses of the theorem, we may apply Theorem~\ref{thm: Ledrappier}
to the Markov chain $\{(X_n,Y_n)\}$ and the matrix function $A'$. This yields a measurable family of probability measures on $\mathbb{P}^{d-1}$,
\[
\{\bar{\eta}_{(x,y)} : (x,y) \in X \times M\},
\]
and a measurable set $U \subset X \times M$ with $(\nu \times m)(U)=1$ such that for every $(x,y) \in U$,
\[
\bar{\eta}_{(x_1,y_1)}
=
Df_x(y)_{*}\,\bar{\eta}_{(x,y)}
\quad
\text{for $P_{(x,y)}$-a.e.\ $(x_1,y_1) \in X \times M$.}
\]
Using the definition of $P_{(x,y)}$, this condition can be rewritten as follows:
for every $(x,y) \in U$,
\begin{equation}
\label{eqn: 1}
\bar{\eta}_{(x_1,f_x(y))}
=
Df_x(y)_{*}\,\bar{\eta}_{(x,y)}
\quad
\text{for $\nu$-a.e.\ $x_1 \in X$.}
\end{equation}
It follows that for each $(x,y) \in U$, the measure $\bar{\eta}_{(x_1,f_x(y))}$ is independent of $x_1$ for $\nu$-almost every $x_1 \in X$.

%\added{New version:}
We now introduce another measurable family $\{\eta_y \colon y \in M\}$ of probability measures on $\Pp^{d-1}$ defined by 
\[
\eta_y(B) = \int_X \bar{\eta}_{(x,y)}(B) \,d \nu(x) 
\]
for every  $y \in M$ and every measurable set $B \subset \Pp^{d-1}$.  Then Property~\eqref{eqn: 1} implies that $Df_x(y)_* \bar{\eta}_{(x,y)} = \eta_{f_x(y)}$ for every $(x,y) \in U$.

Let 
\[
V = \left\{ (x,y) \in X \times M : \bar{\eta}_{(x,y)} = \eta_y \right\}.
\]
This set is measurable, since $(x,y) \mapsto \bar{\eta}_{(x,y)}$ and $(x,y) \mapsto \eta_y$ are measurable maps. 

Next, we show that $(\nu \times m)(V)=1$. For every $x \in X$, let
\[
U_x = \left \{ y \in M : (x,y) \in U\right \}
\]
be the $x$-section of $U$. Since $(\nu \times m)(U)=1$,  Fubini's Theorem implies that there exists $\bar{x} \in X$ such $m(U_{\bar{x}})=1$. By hypothesis, the map $f_{\bar{x}} \colon N_{\bar{x}} \to M$ is an $m$-preserving $C^1$ embedding and $m(N_x)=1$. Therefore, the set $M_{\bar{x}}:= f_{\bar{x}}(U_{\bar{x}} \cap N_{\bar{x}})$ is measurable and $m(M_{\bar{x}})=1$. By the definition of $U_{\bar{x}}$ and Property~\eqref{eqn: 1}, if $y \in M_{\bar{x}}$, then $(x,y) \in V$ for $\nu$-a.e. $x \in X$. Equivalently, 
\[
\nu(V^y)=1 \quad \text{for every } y \in M_{\bar{x}},
\]
where $V^y = \{x \in X : (x,y) \in V\}$ is the $y$-section of $V$. Finally, by Fubini's Theorem,  we obtain
\[
(\nu \times m)(V) = \int_M \nu(V^y) dm(y) \ge  \int_{M_{\bar{x}}} \nu(V^y) dm(y)  =1.
\]

Let $W = U \cap V$. Then $W$ is measurable and satisfies $(\nu \times m)(W)=1$. Take $(x,y) \in W$. Since $(x,y) \in U$, we have $Df_x(y)_* \bar{\eta}_{(x,y)}= \eta_{f_x(y)}$. On the other hand, $(x,y) \in V$ implies that $\bar{\eta}_{(x,y)} = \eta_y$. Combining these identities, we obtain $Df_x(y)_* \eta_y= \eta_{f_x(y)}$.

To complete the proof, consider the probability measure $\eta$ on $PM$ whose disintegration over $m$ is given by $\{\eta_y : y \in M\}$. By the conclusion above, $\eta$ satisfies $(Df_x)_* \eta = \eta$ for $\nu$-a.e. $x \in X$.
%\added{Old version:}
%Let $M_1=f(Z_0)$. We observe that $M_1$ contains a Borel subset of $M$ with full $m$-measure. This follows from the fact that $X \times M$ is a complete separable metric space, that $(\nu \times m)(Z_0)=1$ and that $f_* (\nu \times m) = m$, the latter expressing the stationarity of $m$ for $\{X_n\}$.
%%We observe that the set $M_1$ contains a Borel subset of $M$ with full $m$-measure, which follows from Lemma \ref{lemma:polish}.
%
%We define another measurable family $\{\eta_y \colon y \in M_1\}$ of probability measures on $\Pp^{d-1}$ by 
%\[
%\eta_y(B) = \int_M \bar{\eta}_{(x,y)}(B) \,d \nu(x) 
%\]
%for every measurable set $B \subset \Pp^{d-1}$ and every $y \in M_1$.  According to \eqref{eqn: 1}, for every $y \in M_1$ and $\nu$-a.e. $x \in X$, we have $\bar{\eta}_{(x,y)} = \eta_y$. 
%%Choose one such $x$.  
%Then, using again \eqref{eqn: 1}, we deduce that for every $y \in M_1$ and $\nu$-a.e. $x \in X$,
%\[
%\bar{\eta}_{(x_1,f_{x}(y))} = Df_x(y)_* \eta_y \quad \text{for } \nu \text{-a.e. } x_1 \in X.
%\]
%Integrating both sides with respect to $\nu$ in $x_1$, we obtain that  
%for every $y \in M_1$ and $\nu$-a.e. $x \in X$,
%\begin{equation}
%\label{eq:fiber-invariance}	
%\eta_{f_x(y)} = Df_x(y)_* \eta_y \quad \text{for all } y \in M_1 \text{ and for }  \nu \text{-a.e. } x \in X.
%\end{equation}
%
%To complete the proof, we define the probability $\eta$ on $PM$ as the probability whose disintegration over $m$ is given by $\{\eta_y \colon y \in M_1\}$, and observe that~\eqref{eq:fiber-invariance} implies $\eta = (Df_x)_* \eta$ for $\nu$-a.e. $x \in X$.
\end{proof}

%--- This lemma proves the statement in the paragraph of the previous proof starting with "Let M_1 ..."
%
%\begin{lemma} \label{lemma:polish}
%Let $X$ be a complete separable space with Borel $\sigma$-algebra $\XX$ and Borel probability measure $m$. Let $Y$ be a separable metric space with Borel $\sigma$-algebra $\YY$. Let $T \colon X \to Y$ be a Borel transformation. Consider a set $A \in \XX$. Then there exists a set $B \in \YY$ such that $T_*m(B)=m(A)$.
%\end{lemma}

%\begin{proof}
%Suppose initially that $T$ is continuous. Since $X$ is a complete separable space, there exists an increasing sequence $\{K_n\}$ of compact subsets of $X$ such that $m(K_n) \to m(A)$. Let $B = \bigcup_n T(K_n) \subset T(A)$. Since $T$ is continuous, each $T(K_n)$ is compact, and so $B$ is a Borel set. Moreover, since 
%\[
%T_*m(T(K_n)) = m(T^{-1}(T(K_n))) \ge m(K_n), 
%\]
%we have 
%\[
%m(B) = \lim_n T_*m(T(K_n)) \ge \lim_n m(K_n) = m(A).
%\] 
%Thus, we conclude that there exists a Borel set $B \subset T(A)$ such that $T_*m(B)=m(A)$.

%Now, suppose that $T$ is Borel measurable. By Lusins's Theorem, for every $\epsilon>0$, there exists a compact set $X_\epsilon \subset X$ with  $m(X_\epsilon) \ge 1-\epsilon$ such that $T|_{X_\epsilon} \colon X_\epsilon \to Y$ is continuous, and $m(A) - \epsilon \le m(A \cap X_\epsilon) \le m(A)$. Let $A_\epsilon := A \cap X_\epsilon$. Observe that $X_\epsilon$ is a complete separable space. Then the previous argument for $T$ continuous applied to $T|_{X_\epsilon}$ provides a set $B_\epsilon \in \YY$ such that $B_\epsilon \subset T(A)$ and $ T_*m(B_\epsilon) = m(A_\epsilon)$. Choose $\epsilon_n = 1/n$, and let $B = \bigcup_n B_{\epsilon_n} \in \YY$. Then $B \subset T(A)$ and $T_*m(B_{\epsilon_n}) \to m(A)$.
%\end{proof}


\subsubsection{Random additive perturbations of toral maps}
\label{Consequences}

By Lemma~\ref{le:add-pert}, Theorem~\ref{thm:G} applies to every $\nu$-random additive perturbation $F$ of a toral map $g$. Moreover, for such a system, the condition $\lambda_F^+=\lambda_F^-$ almost everywhere is equivalent to $\lambda_F^+=0$ almost everywhere. Since the projective tangent bundle of $\Tt^d$ is trivial, $P\Tt^d\cong \Tt^d\times\Pp^{d-1}$,
Theorem~\ref{thm:G} immediately yields the following corollary.

\begin{corollary}
\label{cor:H}
Let $F$ be $\nu$-random additive perturbation of a toral map $g$. Suppose that 
\[
\lambda^+_F(\omega,y)=0 \quad \text{for } (\rho_\nu \times m) \text{-a.e. } (\omega,y) \in \Omega \times \Tt^d.
\]
 Then there exists a measurable family of probability measures $\{\eta_y : y \in \Tt^d\}$ on $\Pp^{d-1}$ such that for $m$-a.e. $y \in M$, 
\[
Dg(y)_* \eta_y = \eta_{g(y)+x} \quad \text{for } \nu \text{-a.e. } x \in \Tt^d.
\]
\end{corollary}

%\subsubsection{$\nu$ \added[comment=find a better section title]{absolutely continuous}}


We now examine some consequences of Corollary~\ref{cor:H} under the additional assumption that $\nu$ is absolutely continuous with respect to
$\operatorname{vol}$.

Recall that $\SL^{\pm}(d,\Rr)$ denotes the set of all real $d\times d$ matrices with determinant $\pm1$. Also denote by $\SL(d,\Rr)$ the set of those with determinant
$1$. Let $\mu$ be the push-forward of $m|_N$ under the map $Dg\colon N\to \SL^{\pm}(d,\Rr)$. Let $H\subset \SL^{\pm}(d,\Rr)$ denote the support of $\mu$. For $h\in H$, we use the same symbol to denote both the matrix $h$ and the induced projective transformation of $\Pp^{d-1}$.

%\begin{proposition}
%\label{thm:invariant}
%Let $F$ be random additive perturbation of a map $g$ satisfying Assumption \ref{as: three} with $\nu = \operatorname{vol}_{B_\epsilon(0)}$. Suppose that $\lambda^+_F(\omega,y)=0$ for $\rho_\nu \times m$-a.e. $(\omega,y) \in \Omega \times M$. Then there exists a probability measure $\eta$ on $\Pp^{d-1}$ such that 
%\begin{equation} \label{eqn:8}
%\eta = Dg
%(y)_* \eta \quad \text{for } m \text{-a.e. } y \in M.
%\end{equation}
%Moreover, if we denote by $H$ the closure in $\operatorname{SL}(d,\Rr)$ of the set of matrices $Dg(y)$ satisfying \eqref{eqn:8}, then one of the following must hold: 
%\begin{enumerate}
%\item $H$ is contained in a compact subgroup of $\operatorname{SL}(d,\mathbb{R})$, or 
%\item there exists a nonempty set $L \subset \mathbb{R}^d$ consisting of finitely many proper subspaces such that $h(L)=L$ for every $h\in H$.
%\end{enumerate}
%\end{proposition}

\begin{theorem}
\label{thm:invariant}
Let $F$ be a $\nu$-random additive perturbation of a toral map $g$.
Suppose that $\nu\ll m$ and that $\lambda_F^+(\omega,y)=0$ for $(\rho_\nu\times m)$-a.e. $(\omega,y)\in\Omega\times\Tt^d$. Then there exists a probability measure $\eta$ on $\mathbb P^{d-1}$ such that
\[
h_*\eta=\eta
\quad\text{for all }h\in H.
\]
Moreover, one of the following alternatives holds:
\begin{enumerate}
\item $H$ is contained in a compact subgroup of $\SL^{\pm}(d,\Rr)$,
\item there exists a nonempty finite family $L$ of proper subspaces of $\Rr^d$
such that $h(L)=L$ for every $h\in H$.
\end{enumerate}
\end{theorem}

\begin{proof}
Let $p=d\nu/dm \in L^1(m)$, and set $E=\{x\in\Tt^d:p(x)>0\}$. We have $m(E)>0$. By Corollary~\ref{cor:H}, there exists a measurable
family of probability measures $\{\eta_y:y\in\Tt^d\}$ on $\mathbb P^{d-1}$ and a full $m$-measure set $U\subset N$ such that for every $y\in U$,
\[
Dg(y)_*\eta_y=\eta_{g(y)+x}
\quad\text{for }m\text{-a.e. }x\in E.
\]

Let $V=g(U)$. Since $g|_N$ is a $C^1$ embedding preserving $m$ and $U$ has full $m$-measure, the set $V$ has full $m$-measure. Therefore, for every
$v = g(y) \in V$, we have
\[
Dg(y)_*\eta_y = \eta_{v+x} \quad\text{for }m\text{-a.e. }x\in E.
\]
This implies that the map $w\mapsto\eta_w$ is constant $m$-a.e. on $v+E$. Denote this constant measure by $\zeta_v$.

We now prove that these measures $\zeta_v$ are equal for all $v \in V$. Define
\[
D=\{t\in\Tt^d:m(E\cap(E+t))>0\}.
\]
Since $m(E)>0$, Steinhaus' Theorem~\cite[Corollary 20.17]{hewittross} implies that $D$ contains a neighborhood
$W$ of $0$. If $v,v'\in V$ and $v'-v\in W$, then
\[
m\bigl((v+E)\cap(v'+E)\bigr)
=
m\bigl(E\cap(E+v'-v)\bigr)>0.
\]
Since $w\mapsto\eta_w$ is constant $m$-a.e. on each of the two translates
$v+E$ and $v'+E$, the corresponding constants must coincide.

Choose a symmetric neighborhood $W_0$ of $0$ such that $W_0\subset W$. Since $\Tt^d$ is connected and compact, there exists $N_0\ge 1$ such that
every element of $\Tt^d$ can be written as a sum of $N_0$ elements of $W_0$.
Let $v,v'\in V$, and define $\mathcal C(v,v')$ as the set of tuples
\[
(v_1,\ldots,v_{N_0-1})\in(\Tt^d)^{N_0-1}
\]
such that
\[
v_1-v\in W_0,\quad v_2-v_1\in W_0,\quad \ldots,\quad v'-v_{N_0-1}\in W_0.
\]
Then $\mathcal C(v,v')$ is a nonempty open subset of $(\Tt^d)^{N_0-1}$ and intersects $V^{N_0-1}$, since $V$ has full $m$-measure. Hence any two points $v,v'\in V$ can
be joined by a finite chain
\[
v=v_0,v_1,\ldots,v_{N_0}=v',
\qquad v_i\in V,
\]
with
\[
v_{i+1}-v_i\in W_0
\quad\text{for every }i=0,\ldots,N_0-1.
\]
It follows that $\zeta_v = \zeta_{v'}$ for all $v,v' \in V$. Hence there exists a probability measure $\eta$ on $\mathbb P^{d-1}$ such that for every $v\in V$,
\[
\eta_w=\eta \quad\text{for }m\text{-a.e. }w\in v+E.
\]

Let $A=\{w\in\Tt^d:\eta_w\neq\eta\}$. It is straightforward to see that $A$ is measurable. By the previous conclusion, 
\[
m(A\cap(v+E))=0 \quad \text{for every } v\in V.
\]
By Fubini's Theorem,
\[
0 = \int_V m(A\cap(v+E))\,dm(v) = \int_A m\bigl(V\cap(w-E)\bigr)\,dm(w).
\]
Since $V$ has full $m$-measure, $m(V\cap(w-E))=m(E)>0$ for every $w\in\Tt^d$. Therefore $m(A)=0$, and so
\[
\eta_w=\eta
\quad\text{for }m\text{-a.e. }w\in\Tt^d.
\]

We can then conclude that
\[
Dg(y)_*\eta=\eta
\quad\text{for }m\text{-a.e. }y\in\Tt^d,
\]
or equivalently
\[
A_*\eta=\eta
\quad\text{for }\mu\text{-a.e. }A\in\SL^{\pm}(d,\Rr).
\]
Since the stabilizer of $\eta$ is closed, this identity holds for every
$A\in\supp\mu=H$. Hence
\[
h_*\eta=\eta
\quad\text{for every }h\in H.
\]

The final alternative follows by applying~\cite[Corollary~3.2.2]{zimmer} to the $H$-invariant probability measure $\eta$ on $\mathbb P^{d-1}$ (see also~\cite[Section~7.3]{Viana:2014uo}). Therefore either
$H$ is contained in a compact subgroup of $\SL^{\pm}(d,\Rr)$, or there exists a
nonempty finite family $L$ of proper subspaces of $\Rr^d$ such that
\[
h(L)=L
\quad\text{for every }h\in H.
\]
This proves the proposition.
\end{proof}

\begin{remark} \label{remark:hyp-ell}
The failure of Conditions (1) and (2) in the conclusion of Theorem~\ref{thm:invariant} implies that $\lambda^+_F > 0$ on a set of positive $(\rho_\nu \times m)$-measure. This provides a criterion for the positivity of the maximal Lyapunov exponent of $F$ analogous to Furstenberg's criterion for random matrices~\cite[Section~7.3]{Viana:2014uo}.
\end{remark}


%\begin{remark} \label{remark:dim2}
When $d=2$, the last part of Theorem~\ref{thm:invariant} reduces to the following characterization (see~\cite[Section 6.3]{Viana:2014uo}):  
\begin{enumerate}
    \item $H$ is contained in a compact subgroup of $\SL^{\pm}(2,\mathbb{R})$, or  
    \item there exists a nonempty set $L\subset\mathbb{R}^2$ consisting of one or two $1$-dimensional subspaces such that $h(L)=L$ for every $h\in H$.
\end{enumerate}  
The second case can be further refined into two complementary possibilities:  
\begin{itemize}
    \item $L$ consists of either one or two invariant $1$-dimensional subspaces,
    % meaning that $H$ is contained in either a triangular or a diagonal subgroup.
    \item $L$ consists of two $1$-dimensional subspaces interchanged by some $h \in H$.
\end{itemize}
%\end{remark}

 

%When $d=2$, a sufficient condition for this failure is the existence of elements $h_1, h_2 \in H$ satisfying
%\[
%|\tr(h_1)| > 2 \quad \text{and} \quad 0 < |\tr(h_2)| < 2.
%\]
%The first condition ensures that $h_1$ is hyperbolic, whereas the second yields that $h_2$ is elliptic but not conjugate tto a rotation by angle $\pi/2$.
%%with eigenvalues having nonzero real parts. 
%These properties contradicts Conditions (1) and (2), respectively, and therefore imply that $\lambda^+_F > 0$ a.e. on $\Omega \times \Tt^d$.

The following corollary gives a simple condition ensuring that $\lambda_F^{+}>0$ on a set of positive measure when $d=2$. It will be used in Section~\ref{sec:random_standard} to establish the positivity of $\lambda_F^{+}$ when $F$ is a random additive perturbation of a generalized standard map.

%\added{In the next corollary it is possible to drop the condition $\SL(2)$ by reformulating it as follows: $h_1$ is hyperbolic or parabolic with $h_1 \neq \pm I$ and $h_2$ whose projective action does not have order $2$.}
\begin{corollary}
\label{cor:simple}
Assume that the hypotheses of Theorem~\ref{thm:invariant} hold and that $d=2$. Suppose there exist elements $h_1, h_2 \in H \cap \SL(2,\Rr)$ such that $h_1 \neq \pm I$, $|\tr h_1| \ge 2$ and $0 < |\tr h_2| < 2$. 
%\[
%|\tr(h_1)| > 2 \quad \text{and} \quad 0 < |\tr(h_2)| < 2,
%\]
Then $\lambda^+_F(\omega,y)>0$ on a set of positive $(\rho_\nu \times m)$-measure.
\end{corollary}

\begin{proof}
Since $h_1\neq \pm I$ and $|\operatorname{tr} h_1|\geq 2$, the matrix $h_1$
is either hyperbolic, or parabolic and not diagonalizable. Hence $h_1$ cannot belong to a compact subgroup of $\SL^{\pm}(2,\Rr)$, so Condition~\textup{(1)} in Theorem~\ref{thm:invariant} is violated.

Moreover, since $0<|\operatorname{tr} h_2|<2$, the matrix $h_2$ is elliptic and is not conjugate to a rotation by angle $0$, $\pi/2$, $\pi$ or $3\pi/2$. Therefore the projective action of $h_2$ has neither fixed points nor invariant two-point sets, and Condition~\textup{(2)} is also violated.

Thus neither condition in Theorem~\ref{thm:invariant} holds, and the conclusion follows from that theorem.
\end{proof}

%The conditions $h_1 \neq \pm I$ and $|\tr h_1| \ge 2$ imply that $h_1$ is hyperbolic or parabolic but not diagonalizable, while $0 < |\tr h_2| < 2$ ensures that $h_2$ is elliptic but not conjugate to a rotation by angle $\pi/2$. These properties contradicts Conditions (1) and (2) of Theorem~\ref{thm:invariant}, respectively, and therefore give the desired conclusion.


%The next corollary follows directly from Theorem~\ref{thm:invariant}.

%\begin{corollary}
%\label{cor: J}
%Assume that the hypotheses of Theorem~\ref{thm:invariant} hold. If there exists a %sequence $\{h_n\}_{n \ge 0}$ contained in $H$ and a probability measure
%$\zeta$ on $\mathbb P^{d-1}$ such that $(h_n)_* \eta \longrightarrow \zeta$ as $n %\to + \infty$ in the weak-* topology, then $\zeta=\eta$.
%\end{corollary}

The following corollary will be applied in Section~\ref{sec:billiards}, where we study billiard maps and encounter a sequence of matrices in $H$
with the following degenerating factorization.

\begin{corollary}
\label{cor:K}
Assume that the hypotheses of Theorem~\ref{thm:invariant} hold
and that $d=2$. Suppose that there exists a sequence $\{h_n\}_{n\geq 0}\subset H$ such that
\[
h_n=A_nD_nB_n,
\]
where
\[
A_n=
\begin{pmatrix}
1 & 0\\
0 & \alpha_n
\end{pmatrix},
\qquad
\alpha_n \xrightarrow[n\to\infty]{}0,
\]
\[
D_n=
\begin{pmatrix}
1+a_n & b+b_n\\
c_n & 1+d_n
\end{pmatrix},
\qquad
b\neq0,
\qquad
a_n,b_n,c_n,d_n \xrightarrow[n\to\infty]{}0,
\]
and
\[
B_n=\beta_n
\begin{pmatrix}
\beta_n^{-1} & 0\\
0 & 1
\end{pmatrix},
\qquad
\beta_n \xrightarrow[n\to\infty]{}+\infty.
\]
Then
\[
\eta=\delta_{\hat e},
\]
where $\hat e=[1:0]\in\Pp^1$.
\end{corollary}

\begin{proof}
Let $\widehat{h}_n \colon \mathbb P^1 \to \mathbb P^1$ denote the map induced by the action of $h_n$ on $\mathbb P^1$. It is enough to prove that
\[
(\widehat{h}_n)_*\eta \xrightarrow[n\to\infty]{\mathrm{w}*}\delta_{\hat e},
\]
Indeed, since $h_n\in H$ and $\eta$ is $H$-invariant, it follows that $\eta = \delta_{\hat e}$.

Let $[u:v]\in\Pp^1$. The projective action of $h_n=A_nD_nB_n$ is given by
\[
\widehat h_n([u:v])
=
\left[
\beta_n^{-1}(1+a_n)u+(b+b_n)v:
\alpha_n\beta_n^{-1}c_nu+\alpha_n(1+d_n)v
\right].
\]
If $v\neq0$, then 
\[
\widehat h_n([u:v])\longrightarrow [bv:0]=\hat e.
\]
If $v=0$, then $u\neq0$, and after canceling the common factor
$\beta_n^{-1}u$, we obtain
\[
\widehat h_n([u:0])
=
[1+a_n:\alpha_n c_n]
\longrightarrow \hat e.
\]
Therefore $\widehat h_n(\xi)\to\hat e$ for every $\xi\in\Pp^1$. 

By the Dominated Converge Theorem, for every continuous function $\varphi\colon\Pp^1\to\Rr$,
\[
\int_{\Pp^1}\varphi(\xi)\,d(\widehat h_{n*} \eta)(\xi)
=
\int_{\Pp^1}\varphi(\widehat h_n(\xi))\,d\eta(\xi)
\longrightarrow
\varphi(\hat e).
\]
Hence $\widehat h_{n*} \eta\to\delta_{\hat e}$ in the weak-$*$ topology. 
\end{proof}

%--- Old proof
% \begin{proof}
% By hypothesis, we have
% \[
% h_n
% =
% \beta_n
% \begin{pmatrix}
% \beta_n^{-1}(1+a_n) & b+b_n \\
% \beta_n^{-1}\alpha_n c_n & \alpha_n(1+d_n)
% \end{pmatrix}.
% \]
%
% Let $\widehat{h}_n \colon \mathbb P^1 \to \mathbb P^1$ denote the map induced
% by the action of $h_n$ on $\mathbb P^1$, and let
% $\widehat e \in \mathbb P^1$ be the point with homogeneous coordinates $[1:0]$.
% For any $[u:v] \in \mathbb P^1$, we have
% \begin{align*}
% \widehat{h}_n([u:v])
% &=
% \bigl[
% \beta_n^{-1}(1+a_n)u + (b+b_n)v :
% \beta_n^{-1}\alpha_n c_n u + \alpha_n(1+d_n)v
% \bigr] \\
% &\xrightarrow[n\to\infty]{}
% \begin{cases}
% [bv:0] = \widehat e, & \text{if } v \neq 0, \\
% [u:0] = \widehat e, & \text{if } v=0 \text{ and } u\neq 0.
% \end{cases}
% \end{align*}
% Thus, the sequence of maps $\widehat{h}_n$ converges pointwise to the constant map
% \[
% L \colon \mathbb P^1 \to \mathbb P^1,
% \qquad
% L([u:v]) = \widehat e
% \quad \text{for all } [u:v] \in \mathbb P^1.
% \]
%
% By the Dominated Convergence Theorem, it follows that
% \[
% \widehat{h}_{n*}\eta
% \xrightarrow[n\to\infty]{\;\mathrm{w}^*\;}
% L_{*}\eta = \delta_{\widehat e}.
% \]
% Since $\eta$ is $H$-invariant, we conclude $\eta=\delta_{\widehat e}$.
% \end{proof}

\section{Ergodicity of random additive perturbations of toral maps}
\label{se:ergodicity}

In this section, we study stationary measures for the Markov chains generated by random additive perturbations of toral maps; see Definition~\ref{def:toral_map} for the precise meaning of toral map.

Let $g\colon\Tt^d\to\Tt^d$ be a toral map and let $\nu$ be a probability measure on $\Tt^d$. For each $x\in\Tt^d$, define
\[
f_x\colon\Tt^d\to\Tt^d,
\qquad
f_x(y)=g(y)+x.
\]
Let $(X_n){n\geq 0}$ be an i.i.d.\ sequence of $\Tt^d$-valued random variables with common distribution $\nu$. We consider the Markov chain $(Y_n){n\geq 0}$ on $\Tt^d$ defined by
\[
Y_{n+1}=f_{X_n}(Y_n)=g(Y_n)+X_n.
\]
%\end{equation}
The corresponding transition probabilities are
\[
P_y(E)=\nu\bigl(\{x\in\Tt^d:g(y)+x\in E\}\bigr),
\qquad E\subset \Tt^d \text{ measurable}.
\]
The associated Markov operator is
\[
P\varphi(y)=\int_{\Tt^d}\varphi(g(y)+x)\,d\nu(x)
\]
for every bounded Borel function $\varphi\colon\Tt^d\to\Rr$.

We denote by $P^*$ the dual operator acting on probability measures on $\Tt^d$. A probability measure $\mu$ is called $P$-stationary if $P^*\mu=\mu$. The Markov operator $P$ is called \emph{uniquely ergodic} if it admits a unique stationary probability measure. It is called \emph{strong Feller} if it maps
bounded Borel functions to continuous functions.

We impose additional conditions on the map $g$ and the measure $\nu$. More precisely, we consider the following two cases.

\medskip

\noindent\textbf{Case (i).}
\begin{itemize}
\item $g \colon \Tt^d \to \Tt^d$ is continuous and preserves $m$,
\item $\nu$ is absolutely continuous with respect to $m$.
\end{itemize}

\medskip

\noindent\textbf{Case (ii).}
The natural numebr $d$ is even and write $\Tt^d=\Tt^{d/2}\times\Tt^{d/2}$. Moreover,
\begin{itemize}
\item $g=(g_1,g_2)\colon\Tt^d\to\Tt^{d/2}\times\Tt^{d/2}$ is $C^1$ and
preserves $m$,
\item $\nu=\delta_0\times \chi$, where $\delta_0$ is the Dirac measure at $0\in\Tt^{d/2}$, and $\chi=h\,m_1$; here $m_1$ denotes the normalized Lebesgue measure on the second factor $\Tt^{d/2}$, $h\in L^1(m_1)$, $h\geq 0$, and
\[
\int_{\Tt^{d/2}}h(s)\,dm_1(s)=1,
\]
\item for every $y\in\Tt^d$, the map
\[
H_y\colon\Tt^{d/2}\to\Tt^{d/2},
\qquad
H_y(s)=g_1\bigl(g(y)+(0,s)\bigr),
\]
is a $C^1$ diffeomorphism.
\end{itemize}

These two cases cover the random additive perturbations of billiard maps and generalized standard maps considered later in the paper.

\begin{lemma}
\label{lem:strong-feller-ac}
Suppose that Case (i) holds. Then $P$ is strong Feller.
\end{lemma}

\begin{proof}
Since $\nu\ll m$, there exists $p\in L^1(m)$, $p\geq 0$, such that
\[
d\nu(z)=p(z)\,dm(z).
\]
For $w\in\Tt^d$, let $U_w\colon L^1(m)\to L^1(m)$ be the translation operator
$U_w\psi(z)=\psi(z-w)$. The family $\{U_w:w\in\Tt^d\}$ is strongly continuous on $L^1(m)$.

Let $\varphi\colon\Tt^d\to\Rr$ be bounded and Borel. Using the translation-invariance of $m$, we obtain
\[
\begin{aligned}
P\varphi(y)
&=
\int_{\Tt^d}\varphi(g(y)+z)p(z)\,dm(z) =
\int_{\Tt^d}\varphi(z)p(z-g(y))\,dm(z)  \\
&=
\int_{\Tt^d}\varphi(z)U_{g(y)}p(z)\,dm(z).
\end{aligned}
\]
Therefore, for $y_1,y_2\in\Tt^d$,
\[
\left|P\varphi(y_1)-P\varphi(y_2)\right|
\leq
\|\varphi\|_\infty
\left\|U_{g(y_1)}p-U_{g(y_2)}p\right\|_{L^1}.
\]
Since $g$ is continuous and $w\mapsto U_w p$ is continuous as a map from $\Tt^d$ to $L^1(m)$, the last term converges to $0$ as $y_2\to y_1$. Hence $P\varphi$ is continuous. Thus $P$ is strong Feller.
\end{proof}

\begin{lemma}
\label{lem:strong-feller-singular}
Suppose that Case (ii) holds. Then $Q:=P^2$ is strong Feller.
\end{lemma}

\begin{proof}
Let $\varphi\colon\Tt^d\to\Rr$ be bounded and Borel. Since $\nu=\delta_0\times h\,m_1$, we may write the noise variables as $(0,s)$, with $s\in\Tt^{d/2}$. For $y\in\Tt^d$, define
\[
V_y(s)=g_2\bigl(g(y)+(0,s)\bigr),
\qquad s\in\Tt^{d/2}.
\]
Then
\[
f_{(0,t)}\circ f_{(0,s)}(y)
=
\bigl(H_y(s),V_y(s)+t\bigr),
\]
and therefore
\[
Q\varphi(y)
=
\int_{\Tt^{d/2}}\int_{\Tt^{d/2}}
\varphi\bigl(H_y(s),V_y(s)+t\bigr)h(s)h(t)\,dm_1(s)\,dm_1(t).
\]

Fix $y\in\Tt^d$, and let $\tau_y=H_y^{-1}$. Define
\[
\Phi_y\colon\Tt^{d/2}\times\Tt^{d/2}\to\Tt^d, \qquad \Phi_y(s,t)=\bigl(H_y(s),V_y(s)+t\bigr).
\]
Since $H_y$ is a $C^1$ diffeomorphism, $\Phi_y$ is a $C^1$ diffeomorphism of $\Tt^d$. Its inverse is
\[
\Phi_y^{-1}(z_1,z_2)
=
\bigl(\tau_y(z_1),\,z_2-V_y(\tau_y(z_1))\bigr).
\]
Moreover,
\[
\left|\det D\Phi_y(s,t)\right|
=
\left|\det DH_y(s)\right|.
\]
Thus, by the change of variables $z=\Phi_y(s,t)$, we can write
\[
Q\varphi(y)
=
\int_{\Tt^d}\varphi(z)\alpha_y(z)\,dm(z),
\]
where, for $z=(z_1,z_2)\in\Tt^{d/2}\times\Tt^{d/2}$,
\[
\alpha_y(z_1,z_2)
=
\frac{
h(\tau_y(z_1))
h\bigl(z_2-V_y(\tau_y(z_1))\bigr)
}{
\left|\det DH_y(\tau_y(z_1))\right|
}.
\]

We first prove that $y\mapsto\alpha_y$ is continuous as a map from $\Tt^d$ to $L^1(\Tt^d)$ when $h$ is continuous. Consider
\[
G\colon\Tt^d\times\Tt^{d/2}\to\Tt^d\times\Tt^{d/2},
\qquad
G(y,s)=(y,H_y(s)).
\]
The map $G$ is continuous and bijective. Since the domain is compact and the codomain is Hausdorff, $G^{-1}$ is continuous. Hence
\[
(y,z_1)\mapsto \tau_y(z_1)
\]
is continuous. Moreover, $(y,s)\mapsto V_y(s)$ and $(y,s)\mapsto DH_y(s)$ are continuous. Since $\det DH_y(s)\neq 0$ for every $(y,s)$, compactness implies that
$(y,s)\mapsto \left|\det DH_y(s)\right|^{-1}$ is continuous and bounded. Therefore, if $h$ is continuous, the map $(y,z)\mapsto \alpha_y(z)$ is continuous on $\Tt^d\times\Tt^d$. In particular,
\[
\|\alpha_y-\alpha_{y_0}\|_{L^1}\longrightarrow 0
\qquad\text{as }y\to y_0.
\]

We now pass to the general case $h\in L^1(m_1)$. Choose $h_n\in C(\Tt^{d/2})$ such that $h_n\to h$ in $L^1(m_1)$. Let $\alpha_y^{(n)}$ be the function obtained from the previous formula by replacing $h$ with $h_n$. Equivalently, if $L_y$ denotes the Perron--Frobenius operator associated with the diffeomorphism $\Phi_y$, then
\[
\alpha_y=L_y(h\otimes h),
\qquad
\alpha_y^{(n)}=L_y(h_n\otimes h_n).
\]
Since $L_y$ is an isometry on $L^1(\Tt^d)$, we have
\[
\begin{aligned}
\|\alpha_y-\alpha_y^{(n)}\|_{L^1}
&=
\|L_y(h\otimes h-h_n\otimes h_n)\|_{L^1} \\
&=
\|h\otimes h-h_n\otimes h_n\|_{L^1} \\
&\leq
\|h-h_n\|_{L^1(m_1)}\|h\|_{L^1(m_1)}
+
\|h_n\|_{L^1(m_1)}\|h-h_n\|_{L^1(m_1)}.
\end{aligned}
\]
The right-hand side is independent of $y$ and converges to $0$ as $n\to\infty$. Since $y\mapsto\alpha_y^{(n)}$ is continuous from $\Tt^d$ to $L^1(\Tt^d)$ for each $n$, it follows that
\[
\|\alpha_y-\alpha_{y_0}\|_{L^1}\longrightarrow 0
\qquad\text{as }y\to y_0
\]
also for $h\in L^1(m_1)$.

Consequently,
\[
\left|Q\varphi(y)-Q\varphi(y_0)\right|
\leq
\|\varphi\|_\infty\|\alpha_y-\alpha_{y_0}\|_{L^1}
\longrightarrow 0
\qquad\text{as }y\to y_0.
\]
Thus $Q=P^2$ maps bounded Borel functions to continuous functions. Hence $Q$ is strong Feller.
\end{proof}

\begin{proposition}
\label{prop:unique-ergodicity-additive}
Assume that either Case (i) or Case (ii) holds. Then $m$ is the unique $P$-stationary probability measure. In particular, $m$ is ergodic for $P$.
\end{proposition}

\begin{proof}
First observe that $m$ is $P$-stationary. Indeed, for every bounded Borel function $\varphi\colon\Tt^d\to\Rr$,
\[
\int_{\Tt^d}P\varphi(y)\,dm(y)
=
\int_{\Tt^d}\int_{\Tt^d}\varphi(g(y)+x)\,d\nu(x)\,dm(y).
\]
For each fixed $x\in\Tt^d$, the map $y\mapsto g(y)+x$ preserves $m$, because both $g$ and translations preserve $m$. Hence $P^*m=m$.

Assume first that Case (i) holds. By Lemma~\ref{lem:strong-feller-ac}, $P$ is strong Feller. In addition, $\Tt^d$ is connected, $m$ is $P$-stationary and $\operatorname{supp}m=\Tt^d$. Hence, it follows from~\cite[Part (iii) of Proposition~5.18]{Benaim_2022} that $m$ is the unique $P$-stationary probability measure.

Assume now that Case (ii) holds. By Lemma~\ref{lem:strong-feller-singular}, the operator $Q=P^2$ is strong Feller. Since $m$ is $P$-stationary, it is also $Q$-stationary. Again, $\operatorname{supp}m=\Tt^d$, so~\cite[Part (iii) of Proposition~5.18]{Benaim_2022}, applied to $Q$, implies that $m$ is the unique $Q$-stationary probability measure.

Let $\mu$ be any $P$-stationary probability measure. Then $\mu$ is $Q$-stationary.  By uniqueness of the $Q$-stationary probability measure, $\mu=m$. Therefore $m$ is the unique $P$-stationary probability measure.

Finally, for both cases, uniqueness of the stationary probability measure implies that $m$ is ergodic for $P$.
\end{proof}

\begin{corollary}
\label{cor:skew-product-ergodicity}
Assume that either Case (i) or Case (ii) holds. Then the invariant probability measure $\rho_\nu\times m$ is ergodic for the skew-product $F$.
\end{corollary}

\begin{proof}
By Proposition~\ref{prop:unique-ergodicity-additive}, the measure $m$ is ergodic for the Markov operator $P$. Therefore, by~\cite[Proposition~5.13]{Viana:2014uo}, the invariant probability measure $\rho_\nu\times m$ is ergodic for $F$.
\end{proof}

\begin{remark}
\label{rem:ergodic}
Since $(\rho_\nu\times m)$ is ergodic by Corollary~\ref{cor:skew-product-ergodicity}, the Lyapunov exponents $\lambda_F^-(\omega,y)$ and $\lambda_F^+(\omega,y)$ are constant for $(\rho_\nu\times m)$-almost every $(\omega,y)\in\Omega\times\Tt^d$.
\end{remark}

The unique ergodicity of the Markov operator also gives the following equidistribution result for random orbits.

\begin{proposition}
\label{pr:equidistribution}
Assume that either Case (i) or Case (ii) holds. Then for every $y\in\Tt^d$ and for $\rho_\nu$-almost every $\omega\in\Omega$, the random orbit $\{F_\omega^n(y)\}_{n\geq 0}$ is equidistributed with respect to $m$:
\[
\frac{1}{n}\sum_{k=0}^{n-1}\delta_{F_\omega^k(y)}
\;\xrightarrow[n\to\infty]{\mathrm{w}^*}\;
m \qquad \text{for } \rho_\nu \text{-a.e.} \omega\in\Omega.
\]
In particular, $\rho_\nu$-a.e. random orbit of $y$ is dense in $\Tt^d$.
\end{proposition}

\begin{proof}
%First note that $P$ is Feller. Indeed, if $\varphi\in C(\Tt^d)$, then
%\[
%P\varphi(y)=\int_{\Tt^d}\varphi(g(y)+x)\,d\nu(x)
%\]
%is continuous in $y$, by the continuity of $g$ and dominated convergence.
By Proposition~\ref{prop:unique-ergodicity-additive}, $P$ admits the unique stationary probability measure $m$. Therefore, Breiman's law of large numbers for Markov chains~\cite[Corollary~2.7]{BenoistQuint}, applied to the uniquely ergodic Markov kernel $P$, gives
\[
\frac{1}{n}\sum_{k=0}^{n-1}\delta_{F_\omega^k(y)}
\;\xrightarrow[n\to\infty]{\mathrm{w}^*}\;
m
\]
for $\rho_\nu$-a.e. $\omega\in\Omega$.

Finally, since $m$ assigns positive measure to every nonempty open subset of $\Tt^d$, equidistribution implies that $\rho_\nu$-a.e. random orbit of $y$ is dense in $\Tt^d$.
\end{proof}


%---
\section{Random billiards}
%\section{Vanishing Lyapunov exponents characterize circular billiards}
%\section{Vanishing Lyapunov exponents and integrability in random billiards}
%\section{Convex billiards on surfaces with constant curvature}
\label{sec:billiards}
\subsection{Convex billiards on constant-curvature surfaces}
%\subsection{Billiard map}
%\subsection{Billiard map in coordinates $(s,\theta)$}

In this section, we recall the fundamental definitions and properties of billiards in convex domains on surfaces of constant curvature. For a comprehensive treatment of billiards on general surfaces, we refer to~\cite{DiasCarneiro:2024}, and for the case of surfaces with constant curvature to~\cite{dosSantos:2017}.

Let $S$ denote one of the standard surfaces of constant curvature $K$: the Euclidean plane $\mathbb{E}^2$ ($K=0$), the sphere $\mathbb{S}^2$ ($K=1$) and the hyperbolic plane $\mathbb{H}^2$ ($K=-1$). Let $D \subset S$ be a domain whose boundary $\partial D$ is a $C^2$ simple closed convex curve with positive geodesic curvature (such a curve is called an \emph{oval} in~\cite{dosSantos:2017,DiasCarneiro:2024}). We parametrize $\partial D$ by arc-length $s$, normalized so that the total length of the curve is 1, i.e. $|\partial D|=1$.

A \emph{billiard} with table $D$ is the mechanical system consisting of a point particle that moves inside $D$ along geodesics of $S$ and reflects elastically upon colliding with the boundary $\partial D$, obeying the usual law of reflection: the angles of reflection equals to angle of incidence. We refer to such systems as \emph{convex billiards}.
%Although, this definition applies to a broader class of tables $D$ than those considered in this paper, in what follows, we restrict our attention to tables $D$ as above, and refer to such systems as \emph{convex billiards}.

\subsubsection{Billiard map in coordinates $(s,\theta)$}	
Each collision of the particle with $\partial D$ is described by the pair $(s,\theta)$, where $s$ denotes the arc-length parameter (mod 1) of the point of impact, and $\theta\in[0,\pi]$ is the angle between the incoming trajectory and the positively oriented tangent to $\partial D$ at $s$. 
%Since $\partial D$ is a closed curve of unit length, the arc-length parameter is naturally periodic: two points $s_1$ and $s_2$ with $s_2 - s_1 \in \mathbb{Z}$ represent the same location on the boundary. Therefore, we identify $s$ with a point on the unit circle $S^1$. It follows that 
Hence, the space of all possible collisions is the closed cylinder
\[
Q=S^1 \times [0,\pi]
\] 

The billiard map associated with the table $D$ on a surface $S$ is the map $\phi \colon Q\to Q$ that 
%whose elements are pairs $(s,\theta)$, where $s \in S^1$ indicates the position of the collision along $\partial D$, and $\theta$ is the angle of reflection measured with respect to the positively oriented tangent to $\partial D$ at $s$. The map $\phi$ 
assigns to each collision $(s,\theta) \in Q$ the next collision 
\[
\phi(s,\theta)=(s_1(s,\theta),\theta_1(s,\theta)).
\] 
%In the remainder of the paper, we will refer to such billiard maps as \emph{convex billiards}.

The map $\phi$ for general surfaces of constant curvature retains the same properties of the billiard map for planar convex billiards.

\begin{proposition}[\cite{DiasCarneiro:2024}] \label{prop:properties}
The map $\phi \colon Q \to Q$ satisfies the following properties:
\begin{enumerate} 
\item $\phi$ is a homeomorphism, 
\item $\operatorname{int} Q := S^1 \times (0,\pi)$ and $\partial Q:= S^1 \times \{0,\pi\}$ are $\phi$-invariant sets, 
\item $\phi|_{\operatorname{int} Q} \colon \operatorname{int} Q \to \operatorname{int} Q$ is a twist $C^1$ diffeomorphism,
\item $\phi|_{\partial Q} = \operatorname{id}_{\partial Q}$,
\item $\phi$ preserves the measure $d \bar{m} = \sin \theta \,ds \, d\theta$. 
\end{enumerate}
\end{proposition}

Let $\kappa(s)$ denote the geodesic curvature of $\partial D$ at the point with arc-length parameter $s$. Let $t(s,\theta)$ be the geodesic distance on $S$ between the points of $\partial D$ corresponding to $s$ and $s_1(s,\theta)$. Write 
\[
t = t(s,\theta), \quad \theta_1 = \theta_1(s,\theta), \quad \kappa = \kappa(s), \quad \kappa_1 = \kappa(s_1(s,\theta)).
\] 
Then for every $(s,\theta) \in \operatorname{int} Q$, the derivative $D \phi(s,\theta)$ is given by following expressions, corresponding respectively to the cases $S=\mathbb{E}^2$, $S=\mathbb{S}^2$ and $S=\mathbb{H}^2$~\cite{Chernov:2006we,Katok:1986tg,dosSantos:2017}:
\begin{equation} \label{eq:derivative}
\begin{gathered}
D \phi(s,\theta) = 
\begin{pmatrix}
	\frac{\kappa t -\sin \theta}{\sin \theta_1} & \frac{t}{\sin \theta_1} \\
	\frac{\kappa_1 \kappa t -\kappa_1 \sin \theta -\kappa \sin \theta_1}{\sin \theta_1} & \frac{\kappa_1 t -\sin \theta_1}{\sin \theta_1} 
\end{pmatrix}, \\
D\phi(s,\theta) = 
\begin{pmatrix}
\frac{\kappa \sin t - \cos t \sin\theta}{\sin\theta_1} & \frac{\sin t}{\sin\theta_1}\\
\frac{\sin t (\kappa \kappa_1 - \sin\theta \sin\theta_1) - \cos t (\kappa_1 \sin\theta + \kappa \sin\theta_1)}{\sin\theta_1} & \frac{\kappa_1 \sin t - \cos t \sin\theta_1}{\sin\theta_1}
\end{pmatrix}, \\
D\phi(s,\theta) = 
\begin{pmatrix}
\frac{\kappa \sinh t - \cosh t \sin\theta_1}{\sin\theta_1} & \frac{\sinh t}{\sin\theta_1} \\
\frac{\sinh t (\kappa \kappa_1 + \sin\theta \sin\theta_1) - \cosh t (\kappa_1 \sin\theta + \kappa \sin\theta_1)}{\sin\theta_1} & \frac{\kappa_1 \sinh t - \cosh t \sin\theta_1}{\sin\theta_1}
\end{pmatrix}.
\end{gathered}
\end{equation}

Moreover, in all cases, the derivative $D\phi$ admits the following limits (see~\cite{Katok:1986tg} and \cite[Proposition~3.4]{DiasCarneiro:2024}): for every $s \in S^1$, 
\begin{equation} \label{eq:zeroangle}
\lim_{\theta \to 0^+}  D \phi(s,\theta) = \lim_{\theta \to \pi^-}  D \phi(s,\theta) =
\begin{pmatrix}
	1 & \frac{2}{\kappa(s)} \\
	0 & 1
\end{pmatrix}.
\end{equation}

\subsubsection{Billiard map in coordinates $(s,r)$}
For the purposes of this work, it is convenient to replace the coordinate $\theta$ by $r=-\cos \theta \in [-1,1]$. In coordinates $(s,r)$, the set of all collisions is given by the closed cylinder $V = S^1 \times [-1,1]$, and the billiard map is denoted by $\Phi \colon V \to V$.

The map $\Phi$ inherits the properties of $\phi$ described in Proposition~\ref{prop:properties}. Namely: 1) $\Phi$ is a homeomorphism, 2) the sets $\operatorname{int} V := S^1 \times (-1, 1)$ and $\partial V := S^1 \times \{-1, 1\}$ are invariant under $\Phi$, 3) $\Phi|_{\operatorname{int} V}$ is a twist $C^1$ diffeomorphism, 4) $\Phi_{\partial V} = \operatorname{id}_{\partial V}$, 5) $\Phi$ preserves the measure $ds dr$. 

However, a key difference between the two maps is that $\det D\Phi \equiv 1$ on $\operatorname{int} V$, whereas for $\Phi$, one has  $\det \phi = \sin \theta / \sin \theta_1 \not\equiv 1$ on $\operatorname{int} Q$. This fact is the reason for working with the map $\Phi$ rather than with $\phi$.


%\begin{remark} 
%The property $\det D\Phi \equiv 1$ is the main reason for working with the map $\Phi$ rather than with $\phi$, for which $\det \phi = \sin \theta / \sin \theta_1 \not\equiv 1$. 
%\end{remark}

\subsubsection{Billiard map on the $2$-torus} 
\label{su:torus}
We now introduce a map $T$ on the 2-torus, induced by the billiard map $\Phi$. Random perturbations of $T$ will be studied in Section~\ref{ss:random_billiards}.

Since the restriction of $\Phi$ to $\partial V$ is the identity, we can define an automorphism $T$ of the torus $\Tt^2 = \mathbb{R}/\mathbb{Z} \times \mathbb{R}/2 \mathbb{Z}$ endowed with the flat metric $ds^2+dr^2$ by choosing $R:=[0,1) \times [-1,1)$ as a fundamental domain of $\mathbb{R}/\mathbb{Z} \times \mathbb{R}/2 \mathbb{Z}$ and setting 
\[
T(s,r) = \Phi(s,r) \quad \text{for all } (s,r) \in R.
%[0,1) \times [-1,1).
\] 

The properties of $\Phi$ imply the following properties of $T$: 1) $T\colon\mathbb T^2\to\mathbb T^2$ is a homeomorphism, 2) the sets $\operatorname{int}R=[0,1)\times(-1,1)$ and $\partial R=[0,1)\times\{-1\}$ are $T$-invariant, 3) the restriction $T|_{\operatorname{int}R}\colon\operatorname{int}R\to\operatorname{int}R$ is a $C^1$ diffeomorphism and $\det DT(s,r)=1$ for every $(s,r)\in\operatorname{int}R$, 4) every point of $\partial R$ is fixed by $T$, i.e. 
$T(s,-1)=(s,-1)$ for every $s\in[0,1)$, 5) $T$ preserves the normalized Haar measure $m$ on $\mathbb T^2$, given in the coordinates $(s,r)$ by $dm=c\,ds\,dr$,
where $c>0$ is chosen so that $m(\mathbb T^2)=1$.

\subsection{Circular billiards}
\label{su:circular}

Let $D\subset S$ be a geodesic disk. Since $S$ has constant curvature, the boundary $\partial D$ is necessarily a geodesic circle. In this paper, we restrict our attention to geodesic disks whose boundaries $\partial D$ have positive geodesic curvature.

Under these assumptions, the billiard map $\phi\colon Q\to Q$ in coordinates $(s,\theta)$ associated with the geodesic disk takes the simple form
\[
\phi(s,\theta) = (s+\alpha(\theta),\theta) \quad \text{for all } (s,\theta)\in Q,
\]
for a suitable smooth function $\alpha$ on $\operatorname{int} Q$; see~\cite{Bolotin1992IntegrableBilliards,santos2022break}. Similarly, the corresponding map $T\colon\Tt^2\to\Tt^2$ has the form
\[
T(s,\theta) = (s+\beta(r),r) \quad \text{for all } (s,\theta)\in\Tt^2,
\]
where $\beta(r)=\alpha(-\arccos r)$. Both maps $\phi$ and $T$ are integrable.

%---Bialy's Theorem
%\added[comment=the remaining part of this subsection is no longer needed but add a remark regarding Bialy's theorem after Theorem \ref{thm: 7}]{We now recall a theorem by Bialy} that characterizes circular billiard among all convex billiards on surfaces of constant curvature. Specifically, he proved that circular billiards are the only ones for which every trajectory has no conjugate points~\cite{Bialy:1993ty,Bialy:2013}. Furthermore, he proved that the absence of conjugate points is equivalent to the existence of a measurable monotone invariant sub-bundle -- a concept introduced by Wojtkowski, who also provided an alternative proof of Bialy's theorem in for planar billiards~\cite{Wojtkowski:1994tj}.

%\begin{defn}
%A measurable sub-bundle of $TQ$ 
%\[
%E = \{E(s,\theta) \subset T_{(s,\theta)}Q \colon (s,\theta) \in Q\}
%\] 
%is called \emph{monotone invariant}  if there exist a measurable vector field
%\[
%v(s,\theta)=v_s(s,\theta) \frac{\partial}{\partial s} + v_\theta(s,\theta) \frac{\partial}{\partial \theta} \quad \text{for all } (s,\theta) \in Q
%\] 
%and a positive measurable function $\gamma \colon Q \to (0,+\infty)$ with the following property, for $\bar{m}$-a.e. $(s,\theta) \in \operatorname{int} Q$, 
%\begin{enumerate}
%	\item $v_s(s,\theta)>0$,
%	\item $D\phi(s,\theta)v(s,\theta)=\gamma(s,\theta)v(\phi(s,\theta))$.
%\end{enumerate}
%such that
%\[
%E(s,\theta) = \operatorname{span} \left( v(s,\theta) \right) \quad \text{for all } (s,\theta) \in Q.
%\] 

%\begin{theorem}[{\cite[Theorems~1.1 and 3.1]{Bialy:2013}}]
%\label{thm:Bialy}
%The billiard map $\phi \colon Q \to Q$ admits a monotone invariant sub-bundle of $TQ$ if and only if $D$ is a geodesic disk.
%\end{theorem}

%\begin{remark}
%Bialy used this characterization of the circular tables to give a partial answer to the Birkhoff conjecture, which postulates that the only integrable convex billiards are those with elliptical tables~\cite{Bialy:1993ty}. 
%\end{remark}










%%%%%%%%%%%
%\section{Billiards}
%%\section{Birkhoff billiards}
%\label{sec:Birkhoff billiards}
%\input{birkhoff_billiards.tex}

%%%%%%%%%%%

%\section{Random additive perturbations of the billiard map}
%\section{Random billiards}
%\label{sec:Random additive}
\subsection{Random additive perturbations of billiards}
\label{ss:random_billiards}

Let $T\colon\Tt^2\to\Tt^2$ be the billiard map associated with a convex domain $D$ on a surface of constant curvature. We consider random additive perturbations of $T$ in the framework of Section~\ref{sec:random maps}.

We equip $\Tt^2$ with its standard flat metric. Let $\nu$ be a probability measure on $\Tt^2$, and set $\Omega=(\Tt^2)^{\Nn}$ and $\rho_\nu=\nu^{\Nn}$. For each $x \in X$, define
\[
f\colon \Tt^2\to\Tt^2,
\qquad
f_x(y)=T(y)+x.
\]
The corresponding random dynamical system is given by the skew product
\[
F(\omega,y)=\bigl(\sigma(\omega),f_{\omega_0}(y)\bigr),
\qquad
(\omega,y)\in\Omega\times\Tt^2.
\]

\begin{lemma}
\label{lemma:0}
The billiard map $T\colon\Tt^2\to\Tt^2$ is a toral map in the sense of Definition~\ref{def:toral_map} with $d=2$ and $N:=\operatorname{int}R$. Consequently, for every probability measure $\nu$ on $\Tt^2$, the random dynamical system $F$ is a $\nu$-random additive perturbation of $T$.
\end{lemma}

\begin{proof}
Conditions~(1) and~(2) of Definition~\ref{def:toral_map} follow from the  properties of the map $T$. It remains to verify Condition~(3). Since $DT=D\Phi$ on $N=\operatorname{int}R$ and $m(\partial R)=0$, it suffices to prove that $\log^+\|D\Phi\|$ is integrable with respect to the measure $dsdr$ on $N$. Indeed, since $\Phi$ is a 2-dimensional map, $\det D\Phi \equiv 1$ implies that $\|(D\Phi)^{-1}\|=\|D\Phi\|$, and hence the integrability of $\log^+ \|D\Phi\|$ on $N$ gives the integrability of $\log^+ \|(D\Phi)^{-1}\|$ on $N$. 

Let  $h \colon [0,1) \times (0,\pi) \to N$ be the change of coordinates given by $h(s,\theta)= (s,-\cos \theta)$. Since $\Phi = h \circ \phi \circ h^{-1}$, the chain rule together with
\[
\|Dh(s,\theta)\|\leq 1, \qquad (s,\theta) \in [0,1) \times (0,\pi),
\]
and
\[
\|Dh^{-1}(s,r)\|=\frac{1}{\sqrt{1-r^2}}, \qquad (s,r) \in N,
\]
gives for every $(s,r) \in N$,
\[
\log^+\|D\Phi(s,r)\|
\leq
\log^+\|D\phi(h^{-1}(s,r))\|
-\frac12\log(1-r^2).
\]
The explicit formula for $D\phi$ and the boundedness of the free-flight time function $(s,\theta) \mapsto t(s,\theta)$ yield a constant $C\geq 1$ such that
\[
\|D\phi(h^{-1}(s,r))\|
\leq
\frac{C}{\sqrt{1-r_1(s,r)^2}},
\]
where $\Phi(s,r)=(s_1(s,r),r_1(s,r))$. Therefore for every $(s,r) \in N$,
\[
\log^+\|D\Phi(s,r)\|
\leq
\log C-\frac12\log(1-r^2)
-\frac12\log\bigl(1-r_1(s,r)^2\bigr).
\]
The function $r\mapsto-\log(1-r^2)$ is integrable on $(-1,1)$. Moreover, since $\Phi$ preserves $dsdr$,
\[
\int_N
-\log\bigl(1-r_1(s,r)^2\bigr)\,ds\,dr < \infty, \qquad 
\int_N
-\log(1-r^2)\,ds\,dr<\infty.
\]
Hence $\log^+\|D\Phi\|$ is integrable on $N$. Thus $T$ satisfies Condition~(3). The second assertion of the lemma follows from Lemma~\ref{le:add-pert}.
\end{proof}

%Assumption~\ref{as:one} follows directly from the properties of $T$. In particular, we have $\operatorname{vol}(M\setminus N)=0$.

%To verify Assumption~\ref{as:three}, we need to show that the functions $(x,y) \mapsto \log^+\|Df_x(y)\|$ and $(x,y) \mapsto \log^+\|(Df_x(y))^{-1}\|$ belong to $L^1(\nu \times m)$ for any choice of $\nu$. 

%Since $Df_x = DT$ for every $x \in M$,
%--- Remark on planar billiards
%We remark that the integrability of $\log\|D\phi(h^{-1}(s,r))\|$ could have been deduced directly from the fact that $\phi$ is a $C^1$ diffeomorphism on $Q$~\cite{Katok:1986tg}. However, the proof given above is more general, as it also applies when $\phi$ is only a $C^1$ diffeomorphism on $\operatorname{int} Q$.



%\subsection{}
%Consider the map $T \colon \Tt^2 \to \Tt^2$ derived from the billiard map $\Phi$ at the end of Section~\ref{sec: billiard map} with the measure $m=$ normalized volume. We apply the results of Section~\ref{sec: random maps} to random additive perturbations of $T$. Following the notation of Section~\ref{sec: random maps}, we have $g = T$ and $f_x = \tau_x \circ T = x + T$. It is straightforward to verify that $T$ satisfies the MAIN required properties for $g$ ...

%Let $(\Omega,\Sigma,\rho)$ be the one-sided (two-sided) canonical space of $(N,\mathcal{B}(N))$ with the product measure $\rho=\nu^{\mathbb{N}}$ (or $\rho=\nu^{\mathbb{Z}}$) for some Borel probability $\nu$ on $N$.
%
%Let $F\colon \Omega \times M \to \Omega \times M$ be an additive perturbation of $T$, and let $f$ be the corresponding random additive perturbation with probability space $(\Omega,\Sigma,\rho)$. 

%The map $f$ preserves the probability $\rho \times m$.
%
% and it is straightforward to check that $f$ satisfies Assumptions 1--4 of Section ~\ref{sec: random maps}.
%
%Let $\|\cdot\|$ denote the standard Euclidean norm of $\mathbb{R}^2$. 
%
%$F$ satisfies Assumption~\ref{as: one} (see Remark~\ref{remark: one}). 
%
%
%\begin{lemma}
%\label{lemma: 1}
%$F$ satisfies Assumption~\ref{as: three}. 
%\end{lemma}
%
%\begin{proof} 
%Since $DT_{\omega_0} = DT$ for every $\omega_0 \in N$ and $DT = D\Phi$ on the interior $(0,1) \times (-1,1)$ of the fundamental domain of $\mathbb{T}^2$, it is enough to show that
%$$\log \|D\Phi^{\pm 1}\| \in L^1(\bar{m}),$$
%
%where $\bar{m} = dsdr$ (as a measure of $[0,1] \times [-1,1]$).
% 
%The function $h$ realizing the change of coordinates $(s,r) = h(s,\theta)$ is given by $h(s,\theta)= (s,-\cos \theta)$. 
%
%Let $\phi$ be the billiard map in coordinates $(s,\theta)$. Hence $\Phi = h \circ \phi \circ h^{-1}$ and 
%$$D \Phi(s,r) = D h(\phi(s,\theta)) D \phi(s,\theta) D h^{-1}(s,r).$$
%
%From $\|D \Phi(s,r)\| \le \|Dh(\phi(s,\theta))\| \|D \phi(s,\theta)\| \|Dh^{-1}(s,r)\|$, it follows
%
%\begin{align*} 
%\log \|D \Phi(s,r)\| & \le \log \|Dh(\phi(s,\theta))\| + \log \|D \phi(s,\theta)\| \\ & \quad + \log \|Dh^{-1}(s,r)\|.
%\end{align*}
%
%Since $\|Dh(s,\theta)\|=1$ and $\|Dh^{-1}(s,r)\| = (-\arccos r)' = (1-r^2)^{-1/2}$, we have $\log \|Dh(s,\theta)\|=0$ and 
%
%\begin{align*} 
%\int \int \log \|Dh^{-1}(s,r)\|ds dr & = \int^1_{-1} \log (1-r^2)^{-\frac{1}{2}} dr \\ & = - \frac{1}{2} \int^1_{-1} \log[(1-r)(1+r)]dr \\
%& = -\frac{1}{2} \int^1_{-1} \log(1-r) - \frac{1}{2}\int^1_{-1} \log(1+r) < +\infty.
%\end{align*}
%
%Finally, we observe that under our hypotheses on the domain $D$, the billiard map $\phi$ in coordinate $(s,\theta)$ is a $C^1$ diffeomorphism on $M$ [KS]. This implies $\log \|D\phi\| \in L^1(\bar{m})$.
%
%Putting all the observations above together, we conclude that $\log \|D \Phi\| \in L^1(\bar{m})$.
%
%The proof that $\log \|D \Phi^{-1}\| \in L^1(\bar{m})$ is similar.
%\end{proof}
%
%\begin{remark} 
%For a large class of tables $D$, the billiard map $\phi$ is a $C^1$ diffeomorphism on $\operatorname{int} M$, and this is enough for $\log \|D\phi\| \in L^1(\bar{m})$ (cf. ~\cite{Katok:1986tg,Chernov:2006we}). In the proof above, we could have used this general result instead of using the fact that $\phi$ is a $C^1$ diffeomorphism on the entire $M$. 
%\end{remark}

%Let $F$ be a random additive perturbation of the billiard map $T$. 

\subsection{Vanishing exponents and circular billiards}
\label{sss:abscont}

For notational convenience, we denote the extremal Lyapunov exponents of $F$ by $\lambda^-$ and $\lambda^+$, omitting the subscript $F$. 
%\added[comment=sure about this?]{Additionally, we may simply write a.e. for $\rho_\nu \times m$-a.e. on $\Omega \times M$.}

%Let $\lambda_F^\pm$ be the Lyapunov exponents of $F$.
%In order to simplify the notation, we drop the subscript and write simply $\lambda^\pm$.

We now show that when the billiard table $D$ is a geodesic disk, the Lyapunov exponents $\lambda^-(\omega,y)$ and $\lambda^+(\omega,y)$ vanish for any choice of the probability measure $\nu$.

\begin{lemma}
\label{lemma: 8}
Suppose that $D$ is a geodesic disk. Then for every probability measure $\nu$ on $\Tt^2$, we have 
\[
\lambda^-(\omega,y) = \lambda^+(\omega, y) = 0 \quad \text{for }\rho_\nu \times m \text{-a.e. }(\omega, y) \in \Omega \times \Tt^2.
\]
\end{lemma}

\begin{proof} 
Let $\nu$ be a probability measure on $\Tt^2$. The expression of the billiard map $T$ associated with a geodesic disk $D$ is given in Section~\ref{su:circular}. For every $x\in \Tt^2$, we have
\[
Df_x(s,r)=DT(s,r)
=
\begin{pmatrix}
1 & \beta'(r)\\
0 & 1
\end{pmatrix},
\qquad (s,r)\in \operatorname{int} R.
\]
Hence
\[
Df_x(y)\frac{\partial}{\partial s}
=
\frac{\partial}{\partial s},
\qquad
(x,y)\in \Tt^2\times \operatorname{int} R.
\]
%Since $m(\Tt^2\setminus \operatorname{int} R)=0$ and each random iterate $F_\omega^n$ preserves $m$, the set
%\[
%\left\{(\omega,y)\in \Omega\times \Tt^2 : F_\omega^n(y)\in %\operatorname{int} R\ \text{for all } n\ge 0
%\right\}
%\]
%has full $\rho_\nu\times m$-measure.
It follows that for $\rho_\nu\times m$-a.e. $(\omega,y)$, 
%\[
%D F_\omega^n(y)\frac{\partial}{\partial s}
%=
%\frac{\partial}{\partial s}
%\qquad\text{for all } n\ge 0,
%\]
%and so
\[
\lim_{n\to+\infty}
\frac1n
\log\left\|
D F_\omega^n(y)\frac{\partial}{\partial s}
\right\|
=0,
\]
and hence one Lyapunov exponent vanishes $\rho_\nu\times m$-a.e. Since their sum is zero $\rho_\nu\times m$-a.e., both exponents vanish $\rho_\nu\times m$-a.e., as claimed.
\end{proof}

%It follows that for all $(\omega,y) \in \Omega \times ?$,
%
%\[
%\lambda^+ = \lim_{n \to \infty} \frac{1}{n} \log \left\| Df^{(n)}_{\omega}(y)\right\| = 0 \quad \lambda^- = \lim_{n \to \infty} \frac{1}{n} \log \left\| \left(Df^{(n)}_{\omega}(y)\right)^{-1}\right\| = 0.
%\]

%for all $n \ge 1$,
%\[
% Df^{(n)}_{\omega}(y) \frac{\partial}{\partial s}  =  \frac{\partial}{\partial s} \quad \text{for all } (\omega,y) \in \Omega \times ?.
%\]
%Hence, 
%\[
%\lim_{n \to \infty} \frac{1}{n} \log \left\| Df^{(n)}_{\omega}(y) \frac{\partial}{\partial s} \right\| = 0 \quad \text{for all } (\omega,y) \in \Omega \times ?.
%\]
%By Oseledets' Theorem, we conclude that $\lambda^\pm(\omega, y) = 0$ for 


%\subsection{Vanishing Lyapunov exponents and integrability}
%\subsection{Rigidity in random billiards}
%\subsection{Lyapunov exponents of random billiards}
%We examine two cases separately: 1) $\nu \ll \operatorname{vol}$ and $B_\epsilon(0) \subseteq \supp \nu$ for some $\epsilon>0$, and 2) 
%\[ 
%\nu(E) \ll \int^1_{-1} \chi_E(0,r) dr \quad \text{for all } E \in \mathcal{B}, \quad \{0\} \times [-\epsilon,\epsilon] \subseteq \supp \nu
%\]
%for some $\epsilon>0$.

%\subsubsection{Case $\nu \ll \operatorname{vol}$ and $B_\epsilon(0) \subseteq \supp \nu$}

We now assume that $\nu \ll \operatorname{vol}$, and proceed to prove Theorem~\ref{thm:main}.

For such a $\nu$, all the conclusions of Section~\ref{Consequences} apply to the skew-product $F$. In particular, the probability measure $\rho_\nu \times m$ is ergodic, and the Lyapunov exponents $\lambda^-$ and $\lambda^+$ are constant for $\rho_\nu \times m$-almost every $(\omega, y) \in \Omega \times \Tt^2$. 

Recall that $\phi\colon Q\to Q$ is the billiard map in coordinates $(s,\theta) \in Q$ with $Q = S^1 \times [0,\pi]$.

\begin{proposition}
\label{prop:6}
Suppose $\lambda^+(\omega,y)=0$ for $\rho_\nu \times m$-a.e. $(\omega,y) \in \Omega \times \Tt^2$. Then there exists a strictly positive continuous function $\gamma \colon \operatorname{int} Q \to \mathbb{R}$ such that
\[
D \phi(s,\theta) \frac{\partial}{\partial s} = \gamma(s,\theta) \frac{\partial}{\partial s} \quad \text{for all } (s,\theta) \in \operatorname{int} Q.
\]
\end{proposition}

\begin{proof}
By Lemma~\ref{lemma:0}, Theorem~\ref{thm:invariant} applies to $F$. Moreover, the expression of $D\phi$ in~\eqref{eq:derivative} and the property~\eqref{eq:zeroangle}, together with the definitions of $\Phi$ and $T$, show that $DT$ admits a factorization required by Corollary~\ref{cor:K}, which therefore applies to $F$. Since $DT$ is continuous on $[0,1) \times (-1,1)$, it follows that $\{DT(s,r) \colon (s,r) \in [0,1) \times (-1,1)\}$ is contained in the support of $DT_* m$. Hence, by Theorem~\ref{thm:invariant} and Corollary~\ref{cor:K}, we obtain
\[
DT(s,r)_* \delta_{\widehat{e}} = \delta_{\widehat{e}} \quad \text{for all } (s,r) \in [0,1) \times (-1,1).
\]
 where $\widehat{e} \,$ is the element of $\mathbb{P}^1$ with homogeneous coordinates $[1:0]$. 
 
This implies that $\widehat{e} \,$ is a fixed point of the projective action of $DT(s,r)$ for all $(s,r) \in [0,1) \times (-1,1)$. Equivalently, the subspace $L_s$ spanned by the vector $\partial/\partial s$ is invariant with respect to $DT(s,r)$ for every $(s,r) \in [0,1) \times (-1,1)$.

In view of the definition of $T$, we see that $L_s$ is also invariant under $D\phi(s,\theta)$ for every $(s,\theta) \in \operatorname{int} Q$. Let $\langle \cdot,\cdot \rangle$ denote the Euclidean inner product on $\Rr^2$. Thus, we have
\begin{equation*} 
%\label{eq:invariance}
D\phi(s,\theta) \frac{\partial}{\partial s} = \gamma(s,\theta) \frac{\partial}{\partial s} \quad  \text{for every } (s,\theta) \in \operatorname{int} Q,
\end{equation*}
where $\gamma(s,\theta) := \langle \partial /\partial s,D\phi(s,\theta) \partial /\partial s \rangle$. 
%\added[comment=if $\partial/\partial s$ is an eigenvector then $\gamma$ must be equal to the first entry of the matrix of $D\phi$]{} 
Since $\phi$ is a $C^1$ diffeomorphism on $\operatorname{int} Q$, the function $\gamma$ is well-defined and continuous on $\operatorname{int} Q$. Moreover, $\gamma$ must be either strictly positive or strictly negative, because $\gamma(s,\theta)=0$ for some $(s,\theta) \in \operatorname{int} Q$ would imply $D\phi(s,\theta) \partial/\partial s = 0$, contradicting the invertibility of $D\phi(s,\theta)$. To establsh that $\gamma$ is strictly positive, it suffices to show that for fixed $s$, we have $\lim_{\theta \to 0^+} \gamma(s,\theta)$ exists and is positive. Indeed, by~\eqref{eq:zeroangle}, we have $\lim_{\theta \to 0^+} \gamma(s,\theta) = \lim_{\theta \to 0^+} \langle D\phi(s,\theta) \partial /\partial s,\partial / \partial s \rangle=1$.
%---Old version starting of the last paragraph  
%Since $\phi$ is a $C^1$ diffeomorphism on the entire set $Q$, the function $\gamma$ extends continuously to $Q$. It follows immediately that~\eqref{eq:invariance} extends to the entire $Q$. We claim that $\gamma$ is strictly positive. Suppose by contradiction that $\gamma$ is not strictly positive. Since $\gamma$ is continuous and satisfies $\gamma(s,\theta)=1$ for $\theta\in\{0,\pi\}$ (see the expression of $D\phi$ for $\theta\in\{0,\pi\}$ in Section~\ref{sec: billiard map}), there exists some $(s,\theta)$ with $\gamma(s,\theta)=0$. This would imply $D\phi(s,\theta) \partial/\partial s = 0$, which contradicts the invertibility of $D\phi(s,\theta)$.
\end{proof}

The following is Theorem~\ref{thm:main} written in a form slightly different.

\begin{theorem}
\label{thm:7}
If $D$ is not a geodesic disk, then $\lambda^+(\omega,y)>0$ for $(\rho_\nu\times m)$-a.e. $(\omega,y)\in\Omega\times\Tt^2$. Otherwise, $\lambda^+(\omega,y)=0$ for $(\rho_\nu\times m)$-a.e. $(\omega,y)\in\Omega\times\Tt^2$.
\end{theorem}

\begin{proof}
If $D$ is a geodesic disk, then $\lambda^+ = 0$ for $\rho_\nu \times m$-a.e. $(\omega,y) \in \Omega \times \Tt^2$ by Lemma~\ref{lemma: 8}.

Conversely, assume that $\lambda^+ = 0$ for $\rho_\nu \times m$-a.e. $(\omega,y) \in \Omega \times \Tt^2$. In the remainder of the proof, we first work with the billiard map $\Phi$ in coordinates $(s,r)$, and later return to the map $\phi$. 

Let $E$ be the sub-bundle of $TV$ defined by 
\[
E(s,r) = \operatorname{span} \left(\frac{\partial}{\partial s}\right) \quad \text{for every } (s,r) \in V.
\]
By Proposition~\ref{prop:6}, we have 
\[
D \Phi(s,r) E(s,r) = E(\Phi(s,r)) \quad \text{for all } (s,r) \in \operatorname{int} V. 
\]
Note that the integral curves of the sub-bundle $E$ are the circles $\Gamma_r$, corresponding to the level sets of the function $p(s,r):=r$. Since $E$ is $D \Phi$-invariant, the foliation $\{\Gamma_r\}$ is $\Phi$-invariant.

This implies that $\Phi$ must be a skew-product of the form $\Phi(s,r) = (b(s,r),a(r))$ for some maps $a \colon [-1,1] \to [-1,1]$ and $b \colon V \to S^1$ such that $a$ is a homeomorphism fixing the points $r=-1$ and $r=1$ and a $C^1$ diffeomorphism on $(-1,1)$, whereas $b$ is continuous and $C^1$ on the interior of $V$. Since $\Phi$ preserves $ds dr$ and $p \circ \Phi = a \circ p$, the measure $dr = p_* ds dr$ must be $a$-invariant. This together with the fact that $r=-1$ and $r=1$ are fixed points of $a$ implies that $a$ is the identity. Thus $\Phi(s,r)=(b(s,r),r)$. It follows that each circle $\Gamma_r$ is $\Phi$-invariant.  
	
The same is true for the map $\phi$: it is a skew product of the form $\phi(s,\theta)=(\bar{b}(s,\theta),\theta)$ for some $C^1$ function $\bar{b}$. Note also that in this case $\det D \phi=\sin \theta/\sin \theta_1 = 1$. As a consequence, the entries on the main diagonal of the matrix of $D \phi$ must be equal to 1. Comparing with the expression of $D\phi$ in~\eqref{eq:derivative}, we 
%\added[comment = more details?]
conclude that the curvature $k$ of $\partial D$ is constant. Hence, $D$ is a geodesic disk.

We have proved that, if $\lambda^+=0$ for $(\rho_\nu\times m)$-a.e. $(\omega,y)$, then $D$ is a geodesic disk. Hence, if $D$ is not a geodesic disk, then $\lambda^+>0$ on a set of positive $(\rho_\nu\times m)$-measure. Since $\rho_\nu\times m$ is ergodic for the skew-product $F$ by Corollary~\ref{cor:skew-product-ergodicity}, it follows that $\lambda^+>0$ for $(\rho_\nu\times m)$-a.e. $(\omega,y)$. This completes the proof.
\end{proof}

%--- Old Proof
%Then, by Proposition~\ref{prop: 6}, the sub-bundle $E \subset TQ$ defined by
%\[
%E(s,\theta) = \operatorname{span} \left(\frac{\partial}{\partial s}\right) \quad \text{for every } (s,\theta) \in Q
%\]
%is $D\phi$-invariant.
%
%%Then, by Proposition~\ref{prop: 6}, the sub-bundle $E \subset TQ$ defined by
%%\[
%%E(s,\theta) = \operatorname{span} \left(\frac{\partial}{\partial s}\right) \quad \text{for every } (s,\theta) \in Q
%%\]
%%is monotone invariant. It follows from Theorem~\ref{thm:Bialy} that the domain $D$ must be a geodesic disk.
%%
%%\added{Alternative proof that does not use Bialy's Theorem. The text in red is not needed, it is there to show that this proof can be completed to look like a special case of the Arnold--Liouville Theorem for maps in dim 2. Also there is a remark at the end on "the equal tangent property", that is no longer needed. g is not an exponent!}
% 
%In the rest of the proof, we work with the billiard map $\Phi$ in coordinates $(s,r)$ and we think of $E$ as a sub-bundle of $TV$. The integral curves of the sub-bundle $E$ are the circles $\Gamma_r$ \added[id=g]{corresponding to the level sets of the function $p(s,r)=r$}. Since $E$ is $D \Phi$-invariant, the foliation $\{\Gamma_r\}$ is $\Phi$-invariant, \added[id=g]{i.e. $\Phi(\Gamma_r) = \Gamma_{p(\Phi(s,r))}$ for every $r \in [-1,1]$.}
%
%It follows that $\Phi$ must be a skew-product of the form $\Phi(s,r) = (b(s,r),a(r))$ for some maps $a \colon [-1,1] \to [-1,1]$ and $b \colon V \to S^1$ such that $a$ is a homeomorphism fixing the points $r=-1$ and $r=-1$ and a $C^1$ diffeomorphism on $(-1,1)$, whereas $b$ is continuous and $C^1$ on the interior of $V$. Since $\Phi$ preserves $ds dr$, the measure $dr = p_* ds dr$ must be $a$-invariant. This together the fact that $r=-1$ and $r=1$ are fixed point s of $a$ implies that $a$ is the identity. Thus $\Phi(s,r)=(b(s,r),r)$. It follows that each circle $\Gamma_r$ is $\Phi$-invariant.  \added[id=g]{In other words, $p$ is a first-integral of $\Phi$, i.e. $p \circ \Phi = p$.}
%
%\added[id=g]{Disintegrating $dsdr$ with respect to $dr$, we obtain that $\Phi|_{\Gamma_r}$ preserves $ds$ for every $r$, and we conclude again that $\Phi|_{\Gamma_r}$ is a rigid motion of $S^1$. Since $\Phi$ is orientation preserving, $\Phi|_{\Gamma_r}$ must be a rotation.} 
%	
%The same is true for the map $\phi$: it is a skew product of the form $\phi(s,\theta)=(\bar{b}(s,\theta),\theta)$ for some $C^1$ function $\bar{b}$. Note also that in this case $\det D \phi=\sin \theta/\sin \theta_1 = 1$. As a consequence, the entries on the main diagonal of the matrix of $D \phi$ must be equal to 1. Comparing with the expression of $D\phi$ in~\eqref{eq:derivative}, we conclude that the curvature $k$ of $\partial D$ is constant 	for every surface $S$. Hence, $D$ is a geodesic disk.
%
%\added[id=g]{Remark: to conclude that $D$ is a disk, we could have used also the following fact: since $p$ is a first integral of $\Phi$, the domain $D$ has the \emph{equal tangents property} (see the proof of Lemma 1 for the 3 geometries at the very end of the paper: J. Jer\'onimo-Castro, G. Ruiz-Hern\'andez, S. Tabachikov, The equal tangents property, Advances in Geometry, Volume 14, 3, 2014), which implies that $D$ is a geodesic disk.}
%\end{proof}

%--- Idea for an alternative proof of Theorem \ref{thm: 7}
%\added{Idea for an alternative proof of Theorem \ref{thm: 7} without using Corollary \ref{cor: K}: every convex billiard table $D$ admits at least two 2-periodic points, $y_1$ and $y_2$, corresponding respectively to the diameter and the width of $D$. The derivatives $D\Phi(y_1)$ and $D\Phi(y_2)$ fall into one of the following cases: 1) both parabolic; 2) hyperbolic and parabolic; 3) hyperbolic and elliptic; 4) both hyperbolic. The first two cases can either be ruled out directly or lead to the same conclusion as Corollary \ref{cor: K}. The third case can likely be excluded, possibly after choosing a different $y_2$, as it should be incompatible with the existence of the invariant measure $\eta$; see Remark \ref{remark:hyp-ell}. The fourth case requires further analysis: it should either be ruled out or shown to imply the existence of a unique invariant direction, as in Corollary \ref{cor: K}.
%}

\begin{remark}
The second part of the proof of Theorem~\ref{thm:7} can alternatively be derived from a result of Bialy~\cite{Bialy:1993ty,Bialy:2013}, which characterizes circular billiard tables among all convex tables on surfaces of constant curvature. Specifically, Bialy proved that circular billiards are the only ones for which every trajectory has no conjugate points. Moreover, he showed that the absence of conjugate points is equivalent to the existence of a measurable monotone invariant sub-bundle. By Proposition~\ref{prop:6}, the sub-bundle $E$ introduced in the proof of Theorem~\ref{thm:7} has this property. Thus, Bialy's result applies, and we may conclude that $D$ is a geodesic disk. However, our approach is more direct, since it relies on the specific form of the invariant sub-bundle $E$ derived in Proposition~\ref{prop:6}. For an alternative proof of Bialy's theorem in the planar case, see also~\cite{Wojtkowski:1994tj}.
\end{remark}

%---Old approach: alternative proof based on the isoperimetric inequality

%\added[comment = what follows from here to the end of the section is no longer needed and can be removed]{The second part} of the proof of previous theorem relies Theorem~\ref{thm:Bialy}. However, in the case of planar billiards, the use of that theorem can be replaced by a more direct argument.

%\begin{lemma}
%Suppose that $S=\Ee^2$. Then $\lambda^+(\omega,y)=0$ for $\rho_\nu \times m$-a.e. $(\omega,y) \in \Omega \times M$ if and only if $D$ is a disk.
%\end{lemma}

%\begin{proof}
%The argument is as that of Theorem~\ref{thm: 7}, but we replace the use of Theorem~\ref{thm:Bialy} with a direct computation. From the expression of the matrix of $D\phi$ in~\eqref{eq:derivative}, we see that the conclusion of Proposition~\ref{prop: 6} is equivalent to the identity
%\[
%t(s,\theta) = \frac{\sin \theta}{\kappa(s)} + \frac{\sin \theta_1(s,\theta)}{\kappa(s_1(s,\theta))} \quad \text{for every } (s,\theta) \in \operatorname{int} Q.
%\]
%We now integrate both sides of this equation with respect to the $\phi$-invariant measure $\sin \theta\, ds\, d\theta$. The integral of the left-hand side is given by a well-known formula from integral geometry \cite[Formula~(2.32)]{Chernov:2006we}, 
%\[
%\int_Q t(s,\theta) \sin \theta\, ds\, d\theta = 2\pi\, \operatorname{Area}(D).
%\]
%As for the right-hand side, using the invariance of $\sin \theta ds d\theta$ and the inequality
%\[
%|\partial D|^2 \le 2 \pi \int^{|\partial D|}_0 \frac{1}{\kappa(s)} ds,
%\]
%we obtain
%\begin{gather*}
%\int_Q \left(\frac{\sin \theta}{\kappa(s)} + \frac{\sin \theta_1(s,\theta)}{\kappa(s_1(s,\theta))} \right) \sin \theta ds d\theta  =  2 \int_Q \frac{\sin \theta}{\kappa(s)} \sin \theta ds d\theta \\
% = 2 \int^\pi_0 \sin^2 \theta d\theta \int^{|\partial D|}_0 \frac{1}{\kappa(s)} ds \ge \frac{|\partial D|^2}{2}.
%\end{gather*}
%Therefore, 
%\[
%4 \pi \operatorname{Area}(D) \ge |\partial D|^2.
%\]
%By the isoperimetric inequality, equality must hold. We conclude that $D$ is disk.
%\end{proof}

%\subsection{Magnetic billiards}
%\label{sec:final}

\begin{remark}
We discuss here a possible extension of Theorem~\ref{thm:7} to magnetic billiards. In these systems, a charged point particle moves under the influence of a constant magnetic field. Unlike standard billiards, where the particle travels along geodesics between elastic reflections at the boundary, the motion in magnetic billiards is governed by the Lorentz force, so that trajectories are curves of constant geodesic curvature determined by the field strength;  for instance, see~\cite{Berglund1996MagneticBilliards,Gutkin2001}. Upon collision with the boundary, the particle undergoes specular reflection.

For convex planar tables, Berglund and Kunz~\cite{Berglund1996MagneticBilliards} showed that the magnetic billiard map shares many qualitative features of the corresponding non-magnetic billiard map, with one notable exception: among the two boundary components of the phase space, only the circle $S^{1}\times\{-1\}$ is invariant under the magnetic billiard map, while the other component is not invariant. Nevertheless, the presence of this invariant boundary component suggests that an analog of Theorem~\ref{thm:main} may hold for magnetic billiards with convex planar tables.
\end{remark}

%A second possible extension concerns \added{twist maps}. As already noted, the billiard map $\Phi$ for a billiard with a convex table is a differentiable twist map that preserves the area and fixes each boundary circle  of the cylinder $S^1 \times [-1,1]$. It is straightforward to verify that  all the results in this paper remain valid, up to minor adjustments of the proofs, for every toral map $G \colon \Tt^2 \to \Tt^2$ obtained from a \added[comment=the twist property is not used for $\nu \sim \operatorname{vol}_{B_\epsilon(0)}$; it is for $\nu$ singular]{twist} diffeomorphism $g \colon S^1 \times [a,b] \to S^1 \times [a,b]$ preserving the area and fixing the boundary circles $\Gamma_1 := S^1 \times \{a\}$ and $\Gamma_2 := S^1 \times \{b\}$ by identifying $\Gamma_1$ and $\Gamma_2$. We therefore obtain the following theorem: any random additive perturbation of $G$ has zero Lyapunov exponents if and only if $g$ is the integrable twist map  $g(x,y) = (x+y \bmod 1,y)$.





%--------------------------------------------------

\section{Random generalized standard maps}
%\section{Vanishing Lyapunov exponents characterize the \replaced{constant potential}{integrable} \added{generalized} standard map}
\label{sec:random_standard}
%This section deals with random perturbations of the standard map and contains the proof of Theorem~\ref{thm:standard}, obtained by combining Propositions~\ref{prop:standard-absolute} and~\ref{prop:standard-singular}.

This section deals with random additive perturbations of generalized standard maps. Theorem~\ref{thm:standard} follows from the two results proved below, Proposition~\ref{prop:standard-absolute}, which treats absolutely continuous noise, and Proposition~\ref{prop:standard-singular}, which treats the case of a vertical line noise. 

Write $y=(y_1,y_2) \in \Tt^2$. Recall that a generalized \emph{standard map} is a map $g_V:\Tt^2\to\Tt^2$ defined by
\[
g_V(y_1,y_2)
=
\left(y_1+y_2+V(y_1),\, y_2+V(y_1)\right) \bmod 1,
\]
where $V:\Tt^1\to\Rr$ is a $C^1$ function. Such maps are diffeomorphisms that preserve the normalized Haar measure $m$ on $\Tt^2$. Indeed, 
\[
Dg_V(y_1,y_2)
=
\begin{pmatrix}
1+V'(y_1) & 1\\
V'(y_1) & 1
\end{pmatrix},
\]
and hence $\det Dg_V(y_1,y_2)=1$. If $V\equiv c$ is constant, then $g_V$ is a toral affine skew-product map:
\[
g_V(y_1,y_2)
=
(y_1+y_2+c,\, y_2+c) \bmod 1.
\]
%In particular, for $c=0$ one obtains the integrable map
%\[
%g_0(y_1,y_2)=(y_1+y_2,y_2) \bmod 1.
%\]

Since $g_V$ is a diffeomorphism and preserves $m$, it is a toral map in the sense of Definition~\ref{def:toral_map} with $d=2$ and $N:=\Tt^2$.

In what follows, we use the notation of Subsection~\ref{random additive} with $d=2$. Thus $F \colon \Omega \times M \to \Omega \times M$ denotes the $\nu$-random additive perturbation of $g_V$, and $\lambda^{+}(\omega,y)=\lambda^{+}_F(\omega,y)$ denotes its maximal Lyapunov exponent. Unless otherwise specified, almost everywhere always means with respect to $\rho_\nu \times m$. 

\begin{proposition}
\label{prop:standard-absolute}
Suppose that $\nu \ll m$. If $V$ is constant, then $\lambda^+(\omega,y) = 0$ for $(\rho_\nu \times m)$-a.e. $(\omega,y)$. Otherwise, $\lambda^+(\omega,y) > 0$ for $(\rho_\nu \times m)$-a.e. $(\omega,y)$.
\end{proposition}

\begin{proof}
By Proposition~\ref{prop:unique-ergodicity-additive}, the measure $\rho_\nu \times m$ is ergodic. Therefore there exists a constant $b \ge 0$ such that $\lambda^+ = b$ a.e.

Assume first that $V$ is not constant. Then $V'$ takes both positive and negative values. Consequently, $\operatorname{tr}Dg_V=2+V'$ is larger than $2$ at some points, and belongs to $(0,2)$ at others. Corollary~\ref{cor:simple} implies that $b>0$. 

If $V$ is constant, then
\[
Dg_V \equiv
\begin{pmatrix}
1 & 1 \\
0 & 1
\end{pmatrix}.
\]
It follows that the derivative cocycle has polynomial growth, and therefore $b=0$. This proves the proposition.
\end{proof}

%\added[comment=Note needed since the ergodicity of $\rho_\nu \times m$ follows directly from the conjugation constructed in the proof]{Under Assumption (S), one must also verify the additional non-degeneracy condition. For the maps $g_V$ considered here, this condition is automatic. Indeed, if $g_V(y)=(a,b)$, then
%\[
%H_y(s) = (g_V)_1(a,b+s) = a+b+s+V(a) = C(y)+s,
%\]
%where $C(y)=a+b+V(a)$. Hence $H_y$ is a translation of $\Tt^1$, and therefore a %$C^1$ diffeomorphism.}

\begin{proposition}
\label{prop:standard-singular}
Suppose that $\nu$ is a vertical line measure. If $V$ is constant, then $\lambda^+(\omega,y) = 0$ for $(\rho_\nu \times m)$-a.e. $(\omega,y)$. Otherwise, $\lambda^+(\omega,y) > 0$ for $(\rho_\nu \times m)$-a.e. $(\omega,y)$.
\end{proposition}

\begin{proof}
Recall that $F \colon \Omega\times\Tt^2\to\Omega\times\Tt^2$ is the $\nu$-random additive perturbation of $g_V$. We reduce the proof to the absolutely continuous case by grouping the noise variables $\omega_n$ in pairs.

Recall that $\tau_x$ denotes the translation on $\Tt^2$ by $x \in \Tt^2$, and that $f_x = \tau_x \circ g_V$. A direct computation gives
\[
g_V \circ \tau_{(0,s)} = \tau_{(s,s)} \circ g_V, \quad s \in \Tt^1.
\]
As a consequence, for all $s,t \in \Tt^1$, we have
\begin{equation} \label{eq:two-steps}
f_{(0,t)} \circ f_{(0,s)} = \tau_{(0,t)} \circ g_V \circ \tau_{(0,s)} \circ g_V  = \tau_{(s,s+t)} \circ g^2_V.
\end{equation}

Write $x=(x_1,x_2)$ and $y=(y_1,y_2)$. Define the map $\phi \colon \Tt^2\times\Tt^2\to\Tt^2$ by
\[
\phi(x,y)=(x_1+x_2+y_1,x_2+y_2), \quad x,y \in \Tt^2.
\]
Since $\supp \nu \subset {0}\times \Tt^1$, if $x=(0,s)$ and $y=(0,t)$, then $\phi(x,y)=(s,s+t)$. Therefore, for $(\nu \times \nu)$-a.e. $(\omega_0,\omega_1) \in \Tt^2 \times \Tt^2$, we have
\[
f_{\omega_1} \circ f_{\omega_0} = \tau_{\phi(\omega_0,\omega_1)} \circ g^2_V.
\]

Let $\nu'=\phi_*(\nu\times\nu)$. It is straightforward to check that since $\nu$ is a vertical line measure,  $d\nu' = h(s)h(t-s)dsdt$ for $(s,t) \in \Tt^1 \times \Tt^1$. In particular, $\nu' \ll m$. 

The same observation shows that $\phi$ induces an isomorphism mod $0$ between the probability spaces $(\Tt^2\times\Tt^2,\BB\times\BB,\nu\times\nu)$ and $(\Tt^2,\BB,\nu')$, where $\BB$ is the Borel $\sigma$-algebra of $\Tt^2$.

Define a map $\Phi \colon \Omega \to \Omega$ by
\[
(\Phi(\omega))n =\phi(\omega_{2n},\omega_{2n+1}),
\quad
\omega \in \Omega,\ n \in \Nn.
\]
Let $\rho_{\nu'}=(\nu')^{\Nn}$. By the previous observations on $\phi$, it follows that $\Phi$ is an isomorphism mod $0$ between the Bernoulli shifts $(\Omega,\Sigma,\rho_\nu,\sigma^2)$ and $(\Omega,\Sigma,\rho_{\nu'},\sigma)$, i.e. $\Phi$ is an isomorphism mod $0$ between the probability spaces $(\Omega,\Sigma,\rho_\nu)$ and $(\Omega,\Sigma,\rho_{\nu'})$ such that
\[
\Phi \circ \sigma^2(\omega) = \sigma \circ \Phi(\omega) \quad \text{for } \rho_\nu \text{-a.e. } \omega \in \Omega.
\]

Let $G \colon \Omega \times \Tt^2 \to \Omega \times \Tt^2$ be the $\nu'$-random additive perturbation of $g^2_V$. Define $H \colon \Omega \times \Tt^2 \to \Omega \times \Tt^2$ by
\[
H(\omega,y)=(\Phi(\omega),y), \quad (\omega,y) \in \Omega \times \Tt^2.
\]
Then $H_*(\rho_\nu\times m)=\rho_{\nu'}\times m$, and $H$ conjugates $F^2$ to $G$ mod $0$. Indeed, from~\eqref{eq:two-steps}, it follows that for $\rho_\nu$-a.e. $\omega\in\Omega$ and every $n \ge 1$,
\begin{align}
\label{eq:identity}
G^n_{\Phi(\omega)} & =\tau_{\phi(\omega_{2n-2},\omega_{2n-1})}\circ g_V^{,2}\circ\cdots\circ
\tau_{\phi(\omega_0,\omega_1)}\circ g_V^{,2} \\
&=\tau_{\omega_{2n-1}}\circ g_V\circ\tau_{\omega_{2n-2}}\circ g_V\circ\cdots\circ
\tau_{\omega_1}\circ g_V\circ\tau_{\omega_0}\circ g_V = F^{2n}_\omega.  \notag
\end{align}

Denote by $\lambda^+_G$ and $\lambda^+_F$ the maximal Lyapunov exponents of $G$ and $F$, respectively. For $(\rho_\nu \times m)$-a.e. $(\omega,y) \in \Omega \times \Tt^2$, identity~\eqref{eq:identity} gives
\begin{align*}
\lambda^+_G\left(\Phi(\omega),y\right) & = \lim_{n \to +\infty} \frac{1}{n} \log \|DG^n{\Phi(\omega)}(y)\| \\
& =  \lim_{n \to +\infty} \frac{1}{n} \log  \|DF^{2n}{\omega}(y)\| =  2 \lambda^+_F(\omega,y).
\end{align*}
Since $\Phi_*\rho_\nu=\rho_{\nu'}$, it follows that $\lambda^+_F$ vanishes (is positive) $\rho_\nu\times m$-a.e. if and only if $\lambda^+_G$ vanishes (is positive) $\rho_{\nu'}\times m$-a.e.

We now show that Proposition~\ref{prop:standard-absolute} applies to the system $G$. The measure $\nu'$ is absolutely continuous with respect to $m$. Hence, we only need to verify the trace condition for $g^2_V$. A direct computation yields
\[
\operatorname{tr} Dg^2_V(y_1,y_2) = 2+2V'(y_1) + V'\bigl(y_1+y_2+V(y_1)\bigr)\bigl(2+V'(y_1)\bigr).
\]
Let $\bar y_1\in\Tt^1$ be such that $V'(\bar y_1)=0$. Then
\[
\operatorname{tr} Dg^2_V(\bar y_1,y_2) = 2+2V'\bigl(\bar y_1+y_2+V(\bar y_1)\bigr).
\]
As $y_2$ varies in $\Tt^1$, the argument $\bar y_1+y_2+V(\bar y_1)$ ranges over all of $\Tt^1$. If $V$ is not constant, then $V'$ takes both positive and negative values and vanishes somewhere. Thus the trace of $Dg^2_V$ takes values larger than $2$ and values in $(0,2)$, and the hypotheses of Corollary~\ref{cor:simple} are satisfied.

Therefore the proof of Proposition~\ref{prop:standard-absolute}, applied to $g_V^2$ and $\nu'$, shows that $\lambda_G^+$ is a.e. constant, with zero value if and only if $V$ is constant. The relation between $\lambda_G^+$ and $\lambda_F^+$, then gives the desired conclusion for $\lambda_F^+$.

\end{proof}







%%%%%%%%%%%%%%%%%%%%%%%%%%%%%%%%%%%%%%%%%%%%%%%%%%%%%%%%%%%%

%\section{Random additive perturbations of the outer billiard map}
%\label{sec:Random additive outer}
%\input{random_additive_outer.tex}

%%%%%%%%%%%%%%%%%%%%%%%%%%%%%%%%%%%%%%%%%%%

%\appendix
%\section{Old text}
%\input{appendixA.tex}

%%%%%%%%%%%%%%%%%%%%%%%%%%%%%%%%%%%%%%%%%%%
%\section{Markov chains}

%\input{appendixB.tex}

%%%%%%%%%%%%%%%%%%%%%%%%%%%%%%%%%%%%%%%%%%%
%\section{Lyapunov exponents for linear cocycle over a Markov chain} 

%\input{appendixC.tex}

%%%%%%%%%%%%%%%%%%%%%%%%%%%%%%%%%%%%%%%%%%%%%%%%%%%%%%%%%%%
\section*{Acknowledgements}

%The authors were partially funded by the project ``New trends in Lyapunov exponents’’ PTDC/MAT-PUR/29126/2017.
G. Del Magno acknowledges support from the MIUR Excellence Department Project (CUP I57G22000700001) awarded to the Department of Mathematics, University of Pisa, and from the PRIN Project 2022NTKXCX \textit{``Stochastic properties of dynamical systems''}, funded by the Italian Ministry of University and Research.

JLD and JPG were funded by national funds through FCT – Fundação para a Ciência e a Tecnologia, I.P., in the framework of the unit UID/06522/2025.
%JLD and JPG thank the hospitality of the Department of Mathematics, University of Pisa, Italy, where part of this work was done.

%%%%%%%%%%%%%%%%%%%%%%%%%%%%%%%%%%%%%%%%%%%%%%%%%%%%%%%%%%%%%%%%%%%%%%%%%%%%

\bibliographystyle{abbrv}%{acm}%{plain}
%\addcontentsline{toc}{section}{References}
\bibliography{rfrncs}
\end{document}
